\documentclass[10.5pt,a4paper]{article}
\usepackage[a4paper,top=17mm,bottom=18mm,left=18mm,right=18mm,headsep=6mm,footskip=10mm]{geometry}
\usepackage{fontspec}
\usepackage{kotex}
\IfFontExistsTF{Noto Sans CJK KR}{\setmainfont{Noto Sans CJK KR}\setsansfont{Noto Sans CJK KR}}{}
\usepackage{amsmath,amssymb,mathtools}
\usepackage{booktabs,tabularx,array,longtable,multirow}
\usepackage{enumitem}
\usepackage[most]{tcolorbox}
\usepackage{xcolor}
\usepackage{fancyhdr,lastpage}
\usepackage{hyperref}
\usepackage{tikz}
\usetikzlibrary{arrows.meta,calc,positioning,patterns}
\usepackage{microtype}
\usepackage{setspace}

\definecolor{navy}{HTML}{183153}
\definecolor{blue}{HTML}{2563A6}
\definecolor{sky}{HTML}{EAF3FB}
\definecolor{green}{HTML}{176B52}
\definecolor{mint}{HTML}{EAF7F2}
\definecolor{orange}{HTML}{B45309}
\definecolor{cream}{HTML}{FFF7E8}
\definecolor{red}{HTML}{A33131}
\definecolor{rose}{HTML}{FCEEEE}
\definecolor{graytext}{HTML}{4B5563}
\definecolor{lightgray}{HTML}{F4F6F8}

\hypersetup{colorlinks=true,linkcolor=navy,urlcolor=blue,pdftitle={전반부 공식 PS1--PS4 전 문제 상세풀이},pdfauthor={OpenAI, source-grounded study edition}}
\setlength{\parindent}{0pt}
\setlength{\parskip}{4.5pt}
\setstretch{1.08}
\setlist[itemize]{leftmargin=5mm,itemsep=1.5pt,topsep=2pt}
\setlist[enumerate]{leftmargin=7mm,itemsep=2pt,topsep=2pt}
\renewcommand{\arraystretch}{1.18}
\allowdisplaybreaks

\DeclareMathOperator{\E}{E}
\DeclareMathOperator{\Var}{Var}
\DeclareMathOperator{\Cov}{Cov}
\DeclareMathOperator{\SD}{SD}
\newcommand{\R}{\mathbb{R}}
\newcommand{\dd}{\,\mathrm{d}}
\newcommand{\given}{\,|\,}
\newcommand{\ps}{\mathrm{PS}}
\newcommand{\rbar}{\bar r}
\newcommand{\dbar}{\bar\delta}
\newcommand{\tilden}[1]{\widetilde{#1}}

\pagestyle{fancy}
\fancyhf{}
\fancyhead[L]{\small\color{graytext} Money, Banking and Finance A · 공식 PS1--PS4}
\fancyhead[R]{\small\color{graytext} 1--7강 상세풀이}
\fancyfoot[C]{\small\color{graytext} \thepage\ / \pageref{LastPage}}
\renewcommand{\headrulewidth}{0.3pt}

\newtcolorbox{officialbox}[1][]{enhanced,breakable,colback=sky,colframe=blue,boxrule=0.7pt,arc=2mm,left=3mm,right=3mm,top=2mm,bottom=2mm,title={\bfseries 공식 문제},fonttitle=\normalsize,#1}
\newtcolorbox{strategybox}[1][]{enhanced,breakable,colback=cream,colframe=orange,boxrule=0.7pt,arc=2mm,left=3mm,right=3mm,top=2mm,bottom=2mm,title={\bfseries 풀이 전략},fonttitle=\normalsize,#1}
\newtcolorbox{solutionbox}[1][]{enhanced,breakable,colback=white,colframe=navy,boxrule=0.75pt,arc=2mm,left=3mm,right=3mm,top=2mm,bottom=2mm,title={\bfseries 상세 풀이},fonttitle=\normalsize,#1}
\newtcolorbox{answerbox}[1][]{enhanced,breakable,colback=mint,colframe=green,boxrule=0.8pt,arc=2mm,left=3mm,right=3mm,top=2mm,bottom=2mm,title={\bfseries 최종 답},fonttitle=\normalsize,#1}
\newtcolorbox{warningbox}[1][]{enhanced,breakable,colback=rose,colframe=red,boxrule=0.7pt,arc=2mm,left=3mm,right=3mm,top=2mm,bottom=2mm,title={\bfseries 감점 방지},fonttitle=\normalsize,#1}
\newtcolorbox{sourcebox}[1][]{enhanced,breakable,colback=lightgray,colframe=graytext,boxrule=0.5pt,arc=1.5mm,left=3mm,right=3mm,top=2mm,bottom=2mm,#1}

\newcommand{\chaptertitle}[2]{%
  \clearpage
  \phantomsection
  \addcontentsline{toc}{section}{#1}
  {\color{navy}\LARGE\bfseries #1}\par
  \vspace{1mm}{\color{graytext}\large #2}\par
  \vspace{2mm}\hrule\vspace{4mm}}

\newcommand{\exercise}[2]{%
  \clearpage
  \phantomsection
  \addcontentsline{toc}{subsection}{#1}
  {\color{navy}\Large\bfseries #1}\hfill{\color{graytext}\small #2}\par
  \vspace{1.5mm}\hrule\vspace{3mm}}

\newcommand{\step}[1]{\par\medskip{\color{blue}\bfseries #1}\par}
\newcommand{\finaleq}[1]{\[\boxed{\displaystyle #1}\]}
\newcommand{\smallnote}[1]{{\small\color{graytext}#1}}

\begin{document}

\begin{titlepage}
\thispagestyle{empty}
\vspace*{22mm}
{\color{navy}\Huge\bfseries 전반부 1--7강}\par
\vspace{3mm}
{\color{blue}\fontsize{28}{34}\selectfont\bfseries 공식 PS1--PS4\par 전 문제 상세풀이}\par
\vspace{9mm}
{\Large 교수님이 배포한 공식 Problem Set만 수록}\par
\vspace{3mm}
{\large 임의 제작·변형문제 0개 · 공식 Exercise 23개}\par
\vfill
\begin{tcolorbox}[colback=lightgray,colframe=navy,boxrule=0.8pt,arc=2mm,left=5mm,right=5mm,top=4mm,bottom=4mm]
\textbf{구성 원칙}
\begin{itemize}
\item 문제의 수치·가정·소문항은 교수 배포본 PS1--PS4를 기준으로 보존했다.
\item 공식 해설이 있는 문항은 그 해설의 논리와 기호를 따라 계산 단계를 풀어 썼다.
\item 공식 해설에 ``Discussed in the exercise session''이라고만 적힌 PS3 Exercise 1, PS4 Exercise 1은 수업 녹취와 수업 필기를 근거로 재구성했다.
\item 시험 답안처럼 \textbf{설정식 → FOC 또는 통계량 → 제약조건/시장청산 → 최종 답 → 경제적 직관} 순서로 정리했다.
\end{itemize}
\end{tcolorbox}
\vfill
{\small\color{graytext}Source-grounded Korean study edition · July 2026}\par
\end{titlepage}

\tableofcontents
\clearpage

\section*{자료 범위와 공식성 확인}
\addcontentsline{toc}{section}{자료 범위와 공식성 확인}
\begin{sourcebox}
\begin{tabularx}{\textwidth}{@{}lXl@{}}
\toprule
구분 & 사용한 교수 배포 자료 & 공식 Exercise 수 \\
\midrule
PS1 & MBF\_PS1-2.pdf, MBF\_PS1\_Solution-2.pdf & 4 \\
PS2 & MBF\_PS2\_2024-2.pdf, MBF\_PS2\_Solution\_2024-2.pdf & 5 \\
PS3 & MBF\_PS3-2.pdf, MBF\_PS3\_Solution-2.pdf, PS3\_Ex3,4\_solution.pdf, exercise-session 녹취 & 7 \\
PS4 & MBF\_PS4.pdf, MBF\_PS4\_Solution.pdf, exercise-session 녹취·필기 & 7 \\
\midrule
합계 & 교수 배포 공식 Problem Set & \textbf{23} \\
\bottomrule
\end{tabularx}
\end{sourcebox}

\begin{warningbox}
이 책의 ``문제''는 모두 위 공식 PS의 Exercise다. 설명을 위해 문장을 한국어로 옮겼지만 숫자나 상황을 바꾼 변형문제는 넣지 않았다. 이전의 ``24제'' 자료는 이 공식 문제집과 무관하며 사용하지 않는 것이 맞다.
\end{warningbox}

\section*{시험장에서 공통으로 지킬 답안 순서}
\addcontentsline{toc}{section}{시험장에서 공통으로 지킬 답안 순서}
\begin{enumerate}
\item \textbf{무엇을 선택하는지 먼저 정의}: 소비량, 투자액, 주식 수, 포트폴리오 가중치 중 무엇인지 단위를 적는다.
\item \textbf{제약식 또는 최종부를 먼저 작성}: 예산제약, terminal wealth, 포트폴리오 수익률을 생략하지 않는다.
\item \textbf{공식 전체를 쓴 뒤 특수값을 대입}: 특히 분산식의 공분산 항을 먼저 쓰고, 공분산이 0일 때만 지운다.
\item \textbf{FOC와 제약조건을 함께 사용}: 개인 최적화만 풀고 시장청산을 빼먹지 않는다.
\item \textbf{그림은 축·점·범위를 표시}: 공매도 금지이면 선이 자산점에서 정확히 멈춰야 한다.
\item \textbf{마지막 한 문장 직관}: 위험 증가 → 위험프리미엄 증가 → 가격 하락처럼 방향을 연결한다.
\end{enumerate}

\chaptertitle{PS1 · 일반균형과 Edgeworth Box}{공식 Exercise 1--4 · 소비자 최적화, 시장청산, 계약곡선, 재분배와 가격개입}

\begin{sourcebox}
\textbf{이 장의 공식 출처}: MBF\_PS1-2.pdf, MBF\_PS1\_Solution-2.pdf. 문제의 수치와 효용함수는 그대로 두고, 공식 해설의 계산 사이를 시험 답안 수준으로 보충했다.
\end{sourcebox}

\exercise{PS1 Exercise 1}{공식 문제}
\begin{officialbox}
두 소비자 $i=a,b$와 두 재화 $j=1,2$가 있는 순수교환경제를 생각한다.
\[
(e_{a,1},e_{a,2})=(0,2),\qquad U_a(c_{a,1},c_{a,2})=c_{a,1}+\ln c_{a,2},
\]
\[
(e_{b,1},e_{b,2})=(2,0),\qquad U_b(c_{b,1},c_{b,2})=c_{b,1}^{1/4}c_{b,2}^{3/4}.
\]
재화 $j$의 가격을 $p_j$라 한다.
\begin{enumerate}[label=(\alph*)]
\item 경쟁균형 가격과 배분을 구하라.
\item Edgeworth box에 무차별곡선, 예산선, 계약곡선과 경쟁균형을 표시하라.
\item 시장이 열리기 전에 재화 2만 재분배하여 새 부존량을 $(0,2-x)$와 $(2,x)$로 만들 수 있다. 경쟁균형에서 $c_{a,1}=c_{b,1}=1$을 달성하려면 $x$를 얼마로 해야 하는가? 새 경쟁균형도 구하라.
\item 소비자 $b$의 효용함수가 $\widehat U_b=\ln c_{b,1}+3\ln c_{b,2}$라면 경쟁균형은 어떻게 바뀌는가?
\end{enumerate}
\end{officialbox}

\begin{strategybox}
각 소비자의 \textbf{부존가치가 소득}이다. 개인별 최적화로 수요함수를 만든 뒤, 한 재화의 시장청산으로 상대가격을 구한다. 두 재화 경제에서는 Walras 법칙 때문에 한 시장이 청산되고 예산제약이 만족되면 다른 시장도 청산된다.
\end{strategybox}

\begin{solutionbox}
\step{(a) 소비자 $a$의 수요함수}
$a$의 소득은 $0\cdot p_1+2p_2=2p_2$이다.
\[
\max_{c_{a,1},c_{a,2}}\;c_{a,1}+\ln c_{a,2}
\quad\text{s.t.}\quad p_1c_{a,1}+p_2c_{a,2}=2p_2.
\]
Lagrangian은
\[
\mathcal L_a=c_{a,1}+\ln c_{a,2}
-\lambda\bigl(p_1c_{a,1}+p_2c_{a,2}-2p_2\bigr).
\]
FOC는
\[
1=\lambda p_1,\qquad \frac1{c_{a,2}}=\lambda p_2.
\]
첫 식에서 $\lambda=1/p_1$이므로
\[
\frac1{c_{a,2}}=\frac{p_2}{p_1}
\quad\Longrightarrow\quad c_{a,2}=\frac{p_1}{p_2}.
\]
이를 예산제약에 넣으면
\[
c_{a,1}=\frac{2p_2}{p_1}-1.
\]
따라서
\[
\boxed{c_{a,1}(p)=\frac{2p_2}{p_1}-1,\qquad c_{a,2}(p)=\frac{p_1}{p_2}}.
\]

\step{소비자 $b$의 수요함수}
$b$의 소득은 $2p_1$이다. Cobb--Douglas 효용에서 재화 1과 2의 지출비중은 각각 $1/4$, $3/4$이므로
\[
c_{b,1}=\frac14\frac{2p_1}{p_1}=\frac12,
\qquad
c_{b,2}=\frac34\frac{2p_1}{p_2}=\frac32\frac{p_1}{p_2}.
\]
FOC의 비율로 확인하면
\[
\frac{MU_{b,1}}{MU_{b,2}}
=\frac{c_{b,2}}{3c_{b,1}}=\frac{p_1}{p_2}
\]
도 동일한 결과를 준다.

\step{시장청산과 상대가격}
재화 1의 총부존량은 2이므로
\[
c_{a,1}+c_{b,1}=2.
\]
수요함수를 대입하면
\[
\frac{2p_2}{p_1}-1+\frac12=2
\quad\Longrightarrow\quad
\frac{p_1}{p_2}=\frac45.
\]
가격수준 자체는 임의이므로 $p_1=1$로 정규화하면
\[
(p_1,p_2)=\left(1,\frac54\right).
\]
수요함수에 대입하여
\[
(c_{a,1},c_{a,2})=\left(\frac32,\frac45\right),
\qquad
(c_{b,1},c_{b,2})=\left(\frac12,\frac65\right).
\]
두 소비자의 소비를 합하면 각 재화 모두 2가 되어 시장청산도 확인된다.

\step{(b) 계약곡선}
내부 Pareto 최적배분에서는 두 소비자의 MRS가 같아야 한다.
\[
MRS_a=\frac{MU_{a,1}}{MU_{a,2}}=c_{a,2},
\qquad
MRS_b=\frac{MU_{b,1}}{MU_{b,2}}=\frac{c_{b,2}}{3c_{b,1}}.
\]
자원제약 $c_{b,1}=2-c_{a,1}$, $c_{b,2}=2-c_{a,2}$를 넣으면
\[
c_{a,2}=\frac{2-c_{a,2}}{3(2-c_{a,1})}.
\]
정리하면
\[
\boxed{c_{a,2}=\frac{2}{7-3c_{a,1}}}.
\]
경쟁균형점 $E=(3/2,4/5)$는 이 계약곡선 위에 있고, 예산선은 초기부존점 $I=(0,2)$와 $E$를 지나며 기울기는
\[
-\frac{p_1}{p_2}=-\frac45
\]
이다.

\step{(c) 목표배분에서 역으로 가격과 부존을 찾기}
목표는 $c_{a,1}=c_{b,1}=1$이다. 경쟁균형은 Pareto 효율적이므로 목표배분은 계약곡선 위에 있어야 한다.
\[
c_{a,2}=\frac{2}{7-3\cdot1}=\frac12,
\qquad c_{b,2}=2-\frac12=\frac32.
\]
목표점에서
\[
\frac{p_1}{p_2}=MRS_a=c_{a,2}=\frac12.
\]
$(p_1,p_2)=(1,2)$로 정규화한다. 재분배 후 $a$의 소득은 $(2-x)p_2=2(2-x)$이고, 목표소비의 비용은
\[
p_1\cdot1+p_2\cdot\frac12=1+1=2.
\]
따라서
\[
2(2-x)=2\quad\Longrightarrow\quad \boxed{x=1}.
\]
새 경쟁균형은
\[
(p_1,p_2)=(1,2),\qquad
(c_a,c_b)=\left(\left(1,\frac12\right),\left(1,\frac32\right)\right).
\]

\step{(d) 효용함수의 단조변환}
원래 효용함수에 로그를 취하고 4를 곱하면
\[
4\ln U_b
=4\ln\left(c_{b,1}^{1/4}c_{b,2}^{3/4}\right)
=\ln c_{b,1}+3\ln c_{b,2}=\widehat U_b.
\]
$\widehat U_b$는 $U_b$의 증가변환이므로 동일한 선호와 동일한 무차별곡선을 나타낸다. 따라서 수요와 경쟁균형은 전혀 바뀌지 않는다.
\end{solutionbox}

\begin{answerbox}
\begin{align*}
\text{(a)}\quad &\frac{p_1}{p_2}=\frac45,\quad (p_1,p_2)=\left(1,\frac54\right),\\
&(c_a,c_b)=\left(\left(\frac32,\frac45\right),\left(\frac12,\frac65\right)\right).\\[2mm]
\text{(b)}\quad &c_{a,2}=\frac{2}{7-3c_{a,1}},\qquad \text{예산선 기울기 }-\frac45.\\[2mm]
\text{(c)}\quad &x=1,\quad (p_1,p_2)=(1,2),\\
&(c_a,c_b)=\left(\left(1,\frac12\right),\left(1,\frac32\right)\right).\\[2mm]
\text{(d)}\quad &\text{증가변환이므로 경쟁균형은 변하지 않는다.}
\end{align*}
\end{answerbox}

\begin{warningbox}
가격을 하나의 숫자로만 쓰지 말고 \textbf{정규화 기준}을 밝혀야 한다. 또한 계약곡선은 단순히 두 자산점을 잇는 선이 아니라 $MRS_a=MRS_b$와 자원제약을 함께 풀어야 한다.
\end{warningbox}

\exercise{PS1 Exercise 2}{공식 문제}
\begin{officialbox}
두 소비자의 부존량과 효용함수는
\[
(e_a)=(2,0),\quad U_a=c_{a,1}-e^{-c_{a,2}},
\qquad
(e_b)=(0,2),\quad U_b=c_{b,1}-e^{-c_{b,2}}.
\]
\begin{enumerate}[label=(\alph*)]
\item 경쟁균형을 구하라.
\item Edgeworth box에 무차별곡선, 예산선, 계약곡선과 경쟁균형을 표시하라.
\item 재화 1만 재분배하여 부존량을 $(2-x,0)$와 $(x,2)$로 만들 수 있다. 경쟁균형에서 $c_{a,1}=1$을 달성하려면 $x$는 얼마인가? 새 경쟁균형도 구하라.
\end{enumerate}
\end{officialbox}

\begin{strategybox}
두 소비자의 효용이 준선형이고 동일하다. FOC에서 두 사람 모두 $c_2=\ln(p_1/p_2)$를 선택한다. 재화 2 시장청산이 상대가격을 바로 결정한다.
\end{strategybox}

\begin{solutionbox}
\step{(a) 소비자 $a$}
$a$의 소득은 $2p_1$이다.
\[
\mathcal L_a=c_{a,1}-e^{-c_{a,2}}
-\lambda(p_1c_{a,1}+p_2c_{a,2}-2p_1).
\]
FOC는
\[
1=\lambda p_1,\qquad e^{-c_{a,2}}=\lambda p_2.
\]
따라서
\[
e^{-c_{a,2}}=\frac{p_2}{p_1}
\quad\Longrightarrow\quad
c_{a,2}=\ln\frac{p_1}{p_2}.
\]
예산제약에서
\[
c_{a,1}=2-\frac{p_2}{p_1}\ln\frac{p_1}{p_2}
=2+\frac{p_2}{p_1}\ln\frac{p_2}{p_1}.
\]

\step{소비자 $b$}
$b$의 소득은 $2p_2$이며 동일한 FOC로
\[
c_{b,2}=\ln\frac{p_1}{p_2},
\qquad
c_{b,1}=\frac{p_2}{p_1}\left(2+\ln\frac{p_2}{p_1}\right).
\]

\step{시장청산}
재화 2의 총량은 2이므로
\[
2\ln\frac{p_1}{p_2}=2
\quad\Longrightarrow\quad
\frac{p_1}{p_2}=e.
\]
$p_2=1$로 정규화하면 $(p_1,p_2)=(e,1)$이다. 수요함수에 대입하면
\[
(c_a,c_b)=\left(\left(2-\frac1e,1\right),\left(\frac1e,1\right)\right).
\]

\step{(b) 계약곡선}
\[
MRS_i=\frac{MU_{i,1}}{MU_{i,2}}=e^{c_{i,2}},\qquad i=a,b.
\]
내부 효율성 조건은
\[
e^{c_{a,2}}=e^{c_{b,2}}
\quad\Longrightarrow\quad c_{a,2}=c_{b,2}.
\]
총량이 2이므로
\[
\boxed{c_{a,2}=c_{b,2}=1}.
\]
계약곡선은 Edgeworth box 안의 수평선 $c_{a,2}=1$이고, 예산선 기울기는 $-p_1/p_2=-e$이다.

\step{(c) 재분배}
목표 $c_{a,1}=1$과 계약곡선을 함께 적용하면 목표배분은
\[
(c_a,c_b)=((1,1),(1,1)).
\]
목표점의 MRS가 $e$이므로 가격비율도 $p_1/p_2=e$이다. $(p_1,p_2)=(e,1)$로 두면 $a$의 목표소비 비용은 $e+1$이고 재분배 후 부존가치는 $e(2-x)$다.
\[
e(2-x)=e+1
\quad\Longrightarrow\quad
\boxed{x=1-\frac1e}.
\]
\end{solutionbox}

\begin{answerbox}
\begin{align*}
\text{(a)}\quad &(p_1,p_2)=(e,1),\\
&(c_a,c_b)=\left(\left(2-\frac1e,1\right),\left(\frac1e,1\right)\right).\\
\text{(b)}\quad &\text{계약곡선은 }c_{a,2}=1,\quad\text{예산선 기울기는 }-e.\\
\text{(c)}\quad &x=1-\frac1e,\quad (c_a,c_b)=((1,1),(1,1)),\quad(p_1,p_2)=(e,1).
\end{align*}
\end{answerbox}

\begin{warningbox}
$e^{-c_2}=p_2/p_1$에서 로그를 취할 때 부호를 틀리지 말 것. $c_2=\ln(p_1/p_2)$다. 또한 준선형 효용이라도 예산제약과 비음조건을 확인해야 한다.
\end{warningbox}

\exercise{PS1 Exercise 3}{공식 문제}
\begin{officialbox}
\[
(e_a)=(3,1),\quad U_a=\sqrt{c_{a,1}c_{a,2}},
\qquad
(e_b)=(1,3),\quad U_b=\ln c_{b,1}+\ln c_{b,2}.
\]
\begin{enumerate}[label=(\alph*)]
\item 경쟁균형을 구하고 Edgeworth box에 무차별곡선, 예산선, 계약곡선을 표시하라.
\item 사회계획자가 모든 재화를 $a$에게서 빼앗아 $b$에게 주어 $((0,0),(4,4))$로 만들었다. ``$a$가 아무것도 받지 못하므로 (a)의 배분보다 나쁘다''는 말에 경제학적으로 어떻게 답할 것인가?
\end{enumerate}
\end{officialbox}

\begin{strategybox}
두 효용은 표현만 다를 뿐 모두 두 재화에 동일한 지출비중을 갖는다. 시장청산으로 가격비율이 1이 나온다. (b)는 \textbf{Pareto 효율성과 공정성의 구분}을 묻는다.
\end{strategybox}

\begin{solutionbox}
\step{(a) 개인별 수요}
$a$의 소득은 $3p_1+p_2$다. 효용은 Cobb--Douglas $1/2$--$1/2$이므로
\[
c_{a,1}=\frac{3p_1+p_2}{2p_1}
=\frac32+\frac12\left(\frac{p_1}{p_2}\right)^{-1},
\]
\[
c_{a,2}=\frac{3p_1+p_2}{2p_2}
=\frac12+\frac32\frac{p_1}{p_2}.
\]
$b$의 효용 $\ln c_{b,1}+\ln c_{b,2}$도 두 재화에 절반씩 지출하므로
\[
c_{b,1}=\frac{p_1+3p_2}{2p_1}
=\frac12+\frac32\left(\frac{p_1}{p_2}\right)^{-1},
\]
\[
c_{b,2}=\frac{p_1+3p_2}{2p_2}
=\frac32+\frac12\frac{p_1}{p_2}.
\]
재화 1 시장청산은
\[
c_{a,1}+c_{b,1}=4
\quad\Longrightarrow\quad \frac{p_1}{p_2}=1.
\]
$p_2=1$로 정규화하면
\[
(p_1,p_2)=(1,1),\qquad (c_a,c_b)=((2,2),(2,2)).
\]

\step{계약곡선}
두 소비자의 MRS는
\[
MRS_a=\frac{c_{a,2}}{c_{a,1}},
\qquad
MRS_b=\frac{c_{b,2}}{c_{b,1}}.
\]
자원제약을 넣어
\[
\frac{c_{a,2}}{c_{a,1}}
=\frac{4-c_{a,2}}{4-c_{a,1}}
\quad\Longrightarrow\quad
\boxed{c_{a,1}=c_{a,2}}.
\]
균형점 $(2,2)$는 이 45도 계약곡선 위에 있고 예산선 기울기는 $-1$이다.

\step{(b) 효율성과 공정성}
배분 $((0,0),(4,4))$는 Edgeworth box의 모서리다. 여기서 $a$에게 양의 재화를 조금이라도 주려면 총량이 고정되어 있으므로 $b$의 재화를 줄여야 한다. 두 사람 모두 재화를 더 선호하므로 $b$를 악화시키지 않고 $a$를 개선할 수 없다. 따라서 이 배분은 Pareto 효율적이다.

그러나 Pareto 기준은 두 효율적 배분 중 어느 것이 더 ``공정''한지 순위를 매기지 못한다. (a)의 균형배분과 모서리 배분은 모두 효율적일 수 있고, 공정성·평등성 판단에는 별도의 사회후생함수나 분배규범이 필요하다.
\end{solutionbox}

\begin{answerbox}
\[
(p_1,p_2)=(1,1),\qquad (c_a,c_b)=((2,2),(2,2)),\qquad
\text{계약곡선 }c_{a,1}=c_{a,2}.
\]
$((0,0),(4,4))$도 Pareto 효율적이므로 Pareto 기준만으로 (a)보다 ``나쁘다''고 말할 수 없다. 다만 공정하다는 뜻도 아니다.
\end{answerbox}

\begin{warningbox}
``Pareto efficient''를 ``공정하다'' 또는 ``모두가 행복하다''로 번역하면 오답이다. 효율성은 낭비가 없다는 조건이지 분배의 정의를 보장하지 않는다.
\end{warningbox}

\exercise{PS1 Exercise 4}{공식 문제}
\begin{officialbox}
\[
(e_a)=(3,0),\quad U_a=\sqrt{c_{a,1}}+\frac{c_{a,2}}2,
\qquad
(e_b)=(0,3),\quad U_b=\frac{c_{b,1}}2+\sqrt{c_{b,2}}.
\]
\begin{enumerate}[label=(\alph*)]
\item 경쟁균형을 구하라. 힌트: $x^3+3x^2-3x-1=(x^2+4x+1)(x-1)$.
\item Edgeworth box에 경쟁균형, 무차별곡선, 예산선과 계약곡선을 표시하라.
\item 정부가 재화 1의 가격만 $p_1=3$으로 고정한다. 재화 2 가격은 시장에서 결정된다. 결과적 균형과 (a) 대비 후생변화를 설명하라.
\end{enumerate}
\end{officialbox}

\begin{strategybox}
상대가격 $q=p_1/p_2$를 정의해 수요를 모두 $q$로 표현한다. 힌트의 3차식은 시장청산에서 나온다. (c)는 절대가격이 아니라 상대가격이 실질배분을 결정한다는 점을 묻는다.
\end{strategybox}

\begin{solutionbox}
\step{(a) 소비자 $a$}
$a$의 소득은 $3p_1$이다. Lagrangian의 FOC는
\[
\frac1{2\sqrt{c_{a,1}}}=\lambda p_1,
\qquad
\frac12=\lambda p_2.
\]
두 식의 비율을 취하면
\[
\frac1{\sqrt{c_{a,1}}}=\frac{p_1}{p_2}=q
\quad\Longrightarrow\quad c_{a,1}=q^{-2}.
\]
예산제약에서
\[
c_{a,2}=3q-q^{-1}.
\]

\step{소비자 $b$}
$b$의 소득은 $3p_2$다. FOC는
\[
\frac12=\gamma p_1,
\qquad
\frac1{2\sqrt{c_{b,2}}}=\gamma p_2.
\]
따라서
\[
\sqrt{c_{b,2}}=q,\qquad c_{b,2}=q^2,\qquad
c_{b,1}=\frac3q-q.
\]

\step{시장청산}
재화 1 시장청산은
\[
q^{-2}+\frac3q-q=3.
\]
$q^2$를 곱해 정리하면
\[
q^3+3q^2-3q-1=0.
\]
힌트의 인수분해를 쓰면
\[
(q^2+4q+1)(q-1)=0.
\]
$q^2+4q+1=0$의 두 근은 음수이므로 양의 상대가격은
\[
\boxed{q=\frac{p_1}{p_2}=1}.
\]
$(p_1,p_2)=(1,1)$로 정규화하면
\[
(c_a,c_b)=((1,2),(2,1)).
\]

\step{(b) 계약곡선}
\[
MRS_a=\frac{1}{\sqrt{c_{a,1}}},
\qquad
MRS_b=\sqrt{c_{b,2}}=\sqrt{3-c_{a,2}}.
\]
따라서
\[
\frac1{\sqrt{c_{a,1}}}=\sqrt{3-c_{a,2}}
\quad\Longrightarrow\quad
\boxed{c_{a,2}=3-\frac1{c_{a,1}}}.
\]
예산선은 초기부존점 $(3,0)$과 균형점 $(1,2)$를 지나며 기울기는 $-1$이다.

\step{(c) 정부의 $p_1=3$ 고정}
재화 2의 가격은 자유롭게 조정된다. 시장청산이 요구하는 상대가격은 여전히
\[
\frac{p_1}{p_2}=1.
\]
$p_1=3$이므로 $p_2=3$이 된다. 가격벡터가 $(1,1)$에서 $(3,3)$으로 세 배가 되었지만 상대가격과 각 소비자의 실질소득은 동일하다. 따라서 수요, 배분과 효용은 모두 변하지 않는다.
\end{solutionbox}

\begin{answerbox}
\begin{align*}
\text{(a)}\quad &(p_1,p_2)=(1,1),\quad(c_a,c_b)=((1,2),(2,1)).\\
\text{(b)}\quad &c_{a,2}=3-\frac1{c_{a,1}},\quad\text{예산선 기울기 }-1.\\
\text{(c)}\quad &(p_1,p_2)=(3,3),\quad(c_a,c_b)=((1,2),(2,1)),\\
&\text{두 소비자의 효용은 (a)와 동일하다.}
\end{align*}
\end{answerbox}

\begin{warningbox}
정부가 한 가격의 \textbf{명목수준}만 고정하고 다른 가격이 자유롭게 움직이면, 균형 상대가격을 그대로 복원할 수 있다. $p_1=3$만 보고 상대가격이 3이라고 처리하면 안 된다.
\end{warningbox}

\chaptertitle{PS2 · 위험태도, 확실성등가와 위험프리미엄}{공식 Exercise 1--5 · CARA/CRRA, 기대효용, wealth effect, 확실성등가}

\begin{sourcebox}
\textbf{이 장의 공식 출처}: MBF\_PS2\_2024-2.pdf, MBF\_PS2\_Solution\_2024-2.pdf. 공식 문항의 인물·상황·금액을 그대로 사용했다.
\end{sourcebox}

\exercise{PS2 Exercise 1}{공식 문제}
\begin{officialbox}
Payoff A는 6면체 주사위를 한 번 던져 나온 눈만큼의 달러를 받는 복권이고, Payoff B는 확실하게 3달러를 받는 선택이다.
\begin{enumerate}[label=(\alph*)]
\item 초기부가 0인 Binh의 효용함수는 모든 $W>0$에 대해
\[
-\frac{WU''(W)}{U'(W)}=1
\]
을 만족한다. A와 B 중 무엇을 선택하는가?
\item Guillaume의 효용함수는
\[
-\frac{u''(W)}{u'(W)}=k,\qquad k>0
\]
을 만족한다. 초기부 $Y_0=1$일 때 B를 선택한다는 사실을 알고 있다. 초기부가 충분히 커지면 선택을 바꾸는가? 바꾼다면 임계값을 구하고, 바꾸지 않는다면 그 이유를 설명하라.
\end{enumerate}
\end{officialbox}

\begin{strategybox}
주어진 식은 각각 상대위험회피도(RRA)와 절대위험회피도(ARA)다. 먼저 미분방정식에서 대표 효용함수를 복원한 뒤 기대효용을 비교한다.
\end{strategybox}

\begin{solutionbox}
\step{(a) RRA가 1이면 로그효용}
\[
-\frac{WU''(W)}{U'(W)}=1
\quad\Longrightarrow\quad
\frac{U''(W)}{U'(W)}=-\frac1W.
\]
양변을 적분하면
\[
\ln U'(W)=-\ln W+C
\quad\Longrightarrow\quad U'(W)=\frac{A}{W}.
\]
다시 적분하면 $U(W)=A\ln W+B$다. 양의 affine transformation은 선호를 바꾸지 않으므로 $U(W)=\ln W$로 둔다.

A의 기대효용은
\[
V_A=\frac16\sum_{j=1}^6\ln j
=\frac16\ln(1\cdot2\cdot3\cdot4\cdot5\cdot6)
=\ln\left(720^{1/6}\right).
\]
B의 효용은
\[
V_B=\ln3.
\]
비교하면
\[
V_A<V_B
\quad\Longleftrightarrow\quad
720^{1/6}<3
\quad\Longleftrightarrow\quad
720<3^6=729.
\]
따라서 B를 선택한다.

\step{(b) ARA가 일정하면 CARA 효용}
\[
-\frac{u''(W)}{u'(W)}=k
\quad\Longrightarrow\quad
u'(W)=Ce^{-kW}.
\]
대표 효용함수는
\[
u(W)=-e^{-kW}
\]
로 둘 수 있다. 초기부가 $Y_0$일 때 A의 기대효용은
\[
EU_A(Y_0)=-\frac16\sum_{j=1}^6 e^{-k(Y_0+j)}
=-e^{-kY_0}\left(\frac16\sum_{j=1}^6e^{-kj}\right).
\]
B의 효용은
\[
U_B(Y_0)=-e^{-k(Y_0+3)}=-e^{-kY_0}e^{-3k}.
\]
두 식의 공통 양의 인자 $e^{-kY_0}$는 비교에서 소거된다. 따라서 A와 B의 순위는 $Y_0$와 무관하다. $Y_0=1$에서 B를 선택했다면 모든 초기부에서 B를 선택한다.
\end{solutionbox}

\begin{answerbox}
\begin{enumerate}[label=(\alph*)]
\item $U(W)=\ln W$로 나타낼 수 있고, $720<729$이므로 \textbf{Payoff B}를 선택한다.
\item $u(W)=-e^{-kW}$인 CARA 선호에서는 초기부가 공통인자로 소거된다. 따라서 \textbf{선택을 바꾸지 않으며 임계 초기부도 없다}.
\end{enumerate}
\end{answerbox}

\begin{warningbox}
CARA의 ``wealth effect가 없다''는 말은 위험자산에 투자하는 \textbf{금액} 또는 주어진 복권의 선택이 초기부와 무관하다는 뜻이다. 모든 경제행동이 부와 무관하다는 뜻은 아니다.
\end{warningbox}

\exercise{PS2 Exercise 2}{공식 문제}
\begin{officialbox}
Raluca와 Yuki는 모두
\[
-\frac{WU''(W)}{U'(W)}=\frac12
\]
을 만족하는 선호와 초기부 1달러를 갖는다.
\begin{enumerate}[label=(\alph*)]
\item Argentina와 France가 각각 $1/2$ 확률로 이긴다고 생각한다. ``Argentina 승리 시 Raluca가 Yuki에게서 1달러를 받고, France 승리 시 반대로 지급''하는 내기에 두 사람 모두 동의하는가?
\item Raluca는 Romania 승리확률을 $3/4$, Yuki는 Japan 승리확률을 $3/4$로 믿는다. 동일한 구조의 내기에 두 사람 모두 동의하는가?
\end{enumerate}
\end{officialbox}

\begin{strategybox}
RRA $\gamma=1/2$인 CRRA 효용은 $U(W)=2\sqrt W$로 둘 수 있다. 내기를 수락하면 각 사람의 최종부는 2 또는 0, 거절하면 1이다.
\end{strategybox}

\begin{solutionbox}
\step{효용함수 복원}
CRRA 대표효용은
\[
U(W)=\frac{W^{1-\gamma}}{1-\gamma}.
\]
$\gamma=1/2$를 넣으면
\[
U(W)=2\sqrt W.
\]

\step{(a) 객관적 50--50 내기}
Raluca가 내기를 수락할 때 기대효용은
\[
\frac12U(2)+\frac12U(0)
=\frac12(2\sqrt2)+\frac12(0)=\sqrt2.
\]
거절하면
\[
U(1)=2.
\]
$\sqrt2<2$이므로 Raluca는 거절한다. Yuki에게도 완전히 대칭이므로 Yuki도 거절한다.

\step{(b) 서로 다른 주관적 믿음}
Raluca는 자신이 2달러가 될 확률을 $3/4$로 본다.
\[
EU_R=\frac34U(2)+\frac14U(0)=\frac32\sqrt2.
\]
이를 $U(1)=2$와 비교하면
\[
\frac32\sqrt2\ge2
\quad\Longleftrightarrow\quad
\sqrt2\ge\frac43,
\]
이는 참이다. 따라서 Raluca는 수락한다. Yuki 역시 자신의 관점에서 승리확률을 $3/4$로 믿으므로 같은 계산으로 수락한다.
\end{solutionbox}

\begin{answerbox}
\begin{enumerate}[label=(\alph*)]
\item 둘 다 거절한다. 공정한 평균보존위험을 위험회피자가 선호하지 않기 때문이다.
\item 둘 다 수락한다. 서로 다른 낙관적 믿음 때문에 각자 자신의 승리확률을 $3/4$로 평가한다.
\end{enumerate}
\end{answerbox}

\begin{warningbox}
(b)에서 실제로 양쪽이 동시에 $3/4$ 확률로 이길 수 있다는 뜻이 아니다. \textbf{주관적 확률이 서로 다르기 때문에} 거래가 성립한다는 문제다.
\end{warningbox}

\exercise{PS2 Exercise 3}{공식 문제}
\begin{officialbox}
Raluca의 선호는 모든 $W>0$에 대해
\[
-\frac{U''(W)}{U'(W)}=1
\]
을 만족하고 초기부는 1달러다.
\begin{enumerate}[label=(\alph*)]
\item 두 후보의 승리확률이 각각 $1/2$라고 생각할 때, 이기면 1달러를 받고 지면 1달러를 지급하는 내기를 수락하는가?
\item 다른 선거에서 West가 이기면 1달러를 받고 Trump가 이기면 1달러를 지급하는 내기에 무차별이다. Raluca가 믿는 West의 승리확률을 구하라.
\end{enumerate}
\end{officialbox}

\begin{strategybox}
ARA가 1인 CARA 효용은 $U(W)=-e^{-W}$다. (a)는 공정한 위험의 거절, (b)는 기대효용과 확실한 부 1의 효용을 같게 두는 문제다.
\end{strategybox}

\begin{solutionbox}
\step{(a)}
내기를 수락하면 최종부는 $1/2$ 확률로 2, $1/2$ 확률로 0이다. 기대부는 1로서 거절할 때의 확실한 부 1과 같다. 그러나 $U''<0$인 위험회피자에게
\[
\frac12U(2)+\frac12U(0)<U\left(\frac12\cdot2+\frac12\cdot0\right)=U(1).
\]
따라서 거절한다.

\step{(b)}
West 승리확률을 $\pi$라 하자. 무차별조건은
\[
\pi U(2)+(1-\pi)U(0)=U(1).
\]
$U(W)=-e^{-W}$를 넣으면
\[
-\pi e^{-2}-(1-\pi)=-e^{-1}.
\]
정리하면
\[
\pi(1-e^{-2})=1-e^{-1},
\]
\[
\pi=\frac{1-e^{-1}}{1-e^{-2}}
=\frac{e}{1+e}\approx0.731.
\]
\end{solutionbox}

\begin{answerbox}
\[
\text{(a) 거절},\qquad
\text{(b) }\boxed{\pi=\frac{e}{1+e}\approx73.1\%}.
\]
\end{answerbox}

\begin{warningbox}
(a)에서 ``기대값이 같으므로 무차별''이라고 쓰면 위험중립자의 답이다. 위험회피에서는 Jensen 부등식으로 확실한 1을 더 선호한다.
\end{warningbox}

\exercise{PS2 Exercise 4}{공식 문제}
\begin{officialbox}
게임쇼에서 참가자는 500,000달러를 받고 그만두거나, 네 선택지 중 하나에 답한다. 정답이면 1,000,000달러, 오답이면 32,000달러를 받는다.
\begin{enumerate}[label=(\alph*)]
\item 초기부 0인 Liqian은 RRA가 1이다. 그녀는 현재 500,000달러 제안을 거절하고 답하지만, 확실히 650,000달러를 받을 수 있었다면 그만두었을 것이라고 말한다. 정답확률에 대한 그녀의 믿음을 구하고, 왜 답을 선택하는지 설명하라.
\item 초기부 0인 Dario는 위험선호자이며 네 답이 모두 똑같이 가능하다고 생각한다. 그가 답할지 그만둘지 예측할 수 있는가?
\end{enumerate}
\end{officialbox}

\begin{strategybox}
Liqian은 로그효용을 가지며 650,000달러가 답하기 복권의 확실성등가다. Dario는 단순히 ``위험선호''라는 곡률 정보만 주어졌으므로 효용함수의 정확한 굽힘 정도가 필요하다.
\end{strategybox}

\begin{solutionbox}
\step{(a) Liqian의 주관적 정답확률}
$X_1=1{,}000{,}000$, $X_2=32{,}000$, 확실성등가를 $\widehat K=650{,}000$이라 하자. 정답확률을 $p$라 하면 로그효용의 무차별조건은
\[
p\ln X_1+(1-p)\ln X_2=\ln\widehat K.
\]
이를 풀면
\[
p=\frac{\ln(\widehat K/X_2)}{\ln(X_1/X_2)}
=\frac{\ln(650{,}000/32{,}000)}{\ln(1{,}000{,}000/32{,}000)}
\approx0.875.
\]
답하기 복권의 확실성등가는 650,000달러인데 실제 walk-away 금액은 500,000달러이므로
\[
CE(\text{답하기})=650{,}000>500{,}000.
\]
따라서 답한다.

\step{(b) Dario}
아무 정보가 없으므로 정답확률은 $1/4$이고 오답확률은 $3/4$다. 답하기의 기대금액은
\[
\frac14(1{,}000{,}000)+\frac34(32{,}000)=274{,}000,
\]
500,000달러보다 작다. 위험중립에 가까운 약한 위험선호자라면 그만둘 수 있다. 반면 효용함수가 충분히 강하게 볼록하면 1,000,000달러의 상방을 크게 평가하여 답할 수 있다. $u''(W)>0$만으로는 두 효용을 비교할 수 없다.
\end{solutionbox}

\begin{answerbox}
\[
\text{(a) }p\approx0.875\;(87.5\%),\qquad \text{Liqian은 답한다.}
\]
\[
\text{(b) 판단 불가. 위험선호의 정도, 즉 효용함수의 정확한 곡률이 더 필요하다.}
\]
\end{answerbox}

\begin{warningbox}
위험선호라고 해서 기대금액이 더 낮은 모든 복권을 반드시 고르는 것은 아니다. ``볼록하다''는 방향만 알려 줄 뿐, 두 점 사이의 효용차 크기까지 정하지 않는다.
\end{warningbox}

\exercise{PS2 Exercise 5}{공식 문제}
\begin{officialbox}
Deal or No Deal에서 상자에는 10달러 또는 2달러가 각각 $1/2$ 확률로 들어 있다.
\begin{enumerate}[label=(\alph*)]
\item 초기부 1이고 RRA가 1인 Valérie에게 은행가가 상자를 5달러에 사겠다고 한다. 제안을 수락하는가? 효용곡선으로 설명하라.
\item 초기부가 1보다 크면 결정을 바꾸는가? 바뀌는 임계 초기부를 구하라.
\item 초기부 0이고 RRA가 2인 Pavel에게 은행가가 $B$를 제안한다. Pavel이 수락하는 최소 $B$를 구하라.
\item Pavel의 상자에 대한 위험프리미엄 $\Pi$를 구하라.
\end{enumerate}
\end{officialbox}

\begin{strategybox}
Valérie는 로그효용, Pavel은 $U(W)=-1/W$로 나타낼 수 있다. ``Deal''과 ``No deal''의 최종부를 정확히 비교하고, 최소 제안액은 확실성등가이며 위험프리미엄은 기대값에서 확실성등가를 뺀 값이다.
\end{strategybox}

\begin{solutionbox}
\step{(a) Valérie}
초기부를 $Y=1$, 제안액을 $B=5$, 상자 금액을 $x_1=10$, $x_2=2$라 하자. Deal의 최종부는 $Y+B=6$이고 No deal의 최종부는 11 또는 3이다.
\[
\ln6\;\mathop{\gtrless}_{\text{No deal}}^{\text{Deal}}\;
\frac12\ln11+\frac12\ln3.
\]
로그의 단조성을 이용하면
\[
\ln6>\frac12\ln(33)
\quad\Longleftrightarrow\quad 36>33.
\]
따라서 Deal을 수락한다. 오목한 효용곡선에서 확실한 6의 효용이 11과 3의 기대효용보다 높다.

\step{(b) 초기부의 임계값}
일반적인 $Y$에서 Deal을 선택할 조건은
\[
\ln(Y+B)\ge\frac12\ln(Y+x_1)+\frac12\ln(Y+x_2).
\]
지수화하면
\[
(Y+B)^2\ge(Y+x_1)(Y+x_2).
\]
정리하면
\[
[2B-(x_1+x_2)]Y+B^2-x_1x_2\ge0.
\]
현재 수치를 넣으면
\[
-2Y+5\ge0
\quad\Longleftrightarrow\quad Y\le\frac52.
\]
따라서 $Y=5/2$에서 무차별이고, $Y>5/2$이면 No deal을 택한다. CRRA 효용에서는 초기부가 증가하면 같은 절대금액 복권의 상대적 크기가 작아져 행동이 달라질 수 있다.

\step{(c) Pavel의 최소 제안액}
RRA가 2인 CRRA 효용은
\[
U(W)=\frac{W^{1-2}}{1-2}=-\frac1W
\]
로 둘 수 있다. 초기부가 0이므로 수락조건은
\[
-\frac1B\ge\frac12\left(-\frac1{10}\right)+\frac12\left(-\frac12\right).
\]
오른쪽은 $-3/10$이므로
\[
\frac1B\le\frac3{10}
\quad\Longrightarrow\quad
\boxed{B\ge\frac{10}{3}}.
\]
최소 제안액, 즉 확실성등가는 $\widehat B=10/3$이다.

\step{(d) 위험프리미엄}
상자의 기대금액은
\[
\E[\widetilde x]=\frac12(10)+\frac12(2)=6.
\]
따라서
\[
\Pi=\E[\widetilde x]-CE
=6-\frac{10}{3}=\frac83.
\]
\end{solutionbox}

\begin{answerbox}
\begin{align*}
\text{(a)}\quad &\text{Valérie는 5달러 제안을 수락한다.}\\
\text{(b)}\quad &Y^*=\frac52;\;Y<\frac52\text{이면 Deal, }Y>\frac52\text{이면 No deal}.\\
\text{(c)}\quad &B_{\min}=\frac{10}{3}.\\
\text{(d)}\quad &\Pi=\frac83.
\end{align*}
\end{answerbox}

\begin{warningbox}
위험프리미엄은 $CE-\E[X]$가 아니라 \textbf{$\E[X]-CE$}다. 위험회피자라면 일반적으로 $CE<\E[X]$이므로 위험프리미엄은 양수다.
\end{warningbox}

\chaptertitle{PS3 · 포트폴리오 선택과 균형자산가격}{공식 Exercise 1--7 · 평균--분산, CARA--Normal, 이질적 믿음, 시장청산과 위험프리미엄}

\begin{sourcebox}
\textbf{이 장의 공식 출처}: MBF\_PS3-2.pdf, MBF\_PS3\_Solution-2.pdf, PS3\_Ex3,4\_solution.pdf 및 공식 exercise session 녹취. Exercise 1은 공식 해설 PDF가 ``수업에서 논의''라고만 적고 있으므로 수업시간 풀이를 수식으로 재구성했다.
\end{sourcebox}

\exercise{PS3 Exercise 1}{공식 문제 · 수업 풀이 재구성}
\begin{officialbox}
펀드매니저가 다음 세 자산을 이용한다.
\begin{center}
\begin{tabular}{lccc}
\toprule
자산 & 기대수익률 & 표준편차 & 비고 \\
\midrule
Adidas 주식 & $20\%$ & $4\%$ & 위험자산 \\
Boeing 주식 & $10\%$ & $1.5\%$ & 위험자산 \\
무위험채권 & $2\%$ & $0$ & 무위험자산 \\
\bottomrule
\end{tabular}
\end{center}
Adidas와 Boeing의 상관계수는 $0.2$다. 고객의 초기부는 100,000달러이고
\[
\E[U(W)]=\E[W]-0.0025\Var(W)
\]
를 최대화한다.
\begin{enumerate}[label=(\alph*)]
\item Adidas와 무위험채권만 이용할 때 최적 포트폴리오를 구하라.
\item 세 자산을 모두 이용할 때 최적 포트폴리오를 구하라.
\item 상관계수가 $0.2$에서 $-0.2$로 하락하면 포트폴리오를 어떻게 바꾸는가? 경제적 직관을 설명하라.
\end{enumerate}
\end{officialbox}

\begin{strategybox}
투자액을 각각 $a,b$라 하고 채권투자액을 $W_0-a-b$라 둔다. 먼저 최종부를 쓰고, 기대값과 분산을 구한 뒤 $a,b$로 미분한다. 수업에서 강조한 대로 처음부터 숫자를 난입하기보다 일반식을 먼저 세운다.
\end{strategybox}

\begin{solutionbox}
\step{공통 설정}
순수익률을 $\widetilde r_A,\widetilde r_B$, 무위험 순수익률을 $r_f$라 하자.
\[
\widetilde W_1=(1+\widetilde r_A)a+(1+\widetilde r_B)b+(1+r_f)(W_0-a-b).
\]
따라서
\[
\E[\widetilde W_1]=(1+r_f)W_0+(\bar r_A-r_f)a+(\bar r_B-r_f)b,
\]
\[
\Var(\widetilde W_1)=a^2\sigma_A^2+b^2\sigma_B^2+2ab\sigma_{AB}.
\]
문제의 효용은 $\E[W]-(\theta/2)\Var(W)$ 형태이며 $\theta/2=0.0025$, 즉 $\theta=0.005$다.

\step{(a) Adidas와 채권만}
$b=0$으로 두면 목적함수에서 $a$와 관련된 부분은
\[
(0.20-0.02)a-0.0025(0.04^2)a^2.
\]
FOC는
\[
0.18-2(0.0025)(0.04^2)a=0.
\]
따라서
\[
a^*=\frac{0.18}{0.005\times0.04^2}=22{,}500.
\]
채권에는
\[
100{,}000-22{,}500=77{,}500
\]
달러를 투자한다.

\step{(b) 세 자산 모두}
공분산은
\[
\sigma_{AB}=\rho_{AB}\sigma_A\sigma_B
=0.2(0.04)(0.015)=0.00012.
\]
FOC는
\[
0.18=0.005\bigl(0.0016a+0.00012b\bigr),
\]
\[
0.08=0.005\bigl(0.00012a+0.000225b\bigr).
\]
즉
\[
\begin{pmatrix}
0.0016&0.00012\\
0.00012&0.000225
\end{pmatrix}
\begin{pmatrix}a\\b\end{pmatrix}
=
\begin{pmatrix}36\\16\end{pmatrix}.
\]
행렬식은
\[
D=0.0016(0.000225)-(0.00012)^2=3.456\times10^{-7}.
\]
Cramer 법칙 또는 연립방정식으로
\[
a^*=\frac{36(0.000225)-16(0.00012)}{D}
\approx17{,}881.94,
\]
\[
b^*=\frac{16(0.0016)-36(0.00012)}{D}
\approx61{,}574.07.
\]
따라서 채권투자액은
\[
100{,}000-a^*-b^*\approx20{,}543.98.
\]

\step{(c) 상관계수 $-0.2$}
새 공분산은 $-0.00012$다. FOC의 공분산 항 부호만 바뀐다.
\[
\begin{pmatrix}
0.0016&-0.00012\\
-0.00012&0.000225
\end{pmatrix}
\begin{pmatrix}a\\b\end{pmatrix}
=
\begin{pmatrix}36\\16\end{pmatrix}.
\]
행렬식은 공분산의 제곱이 들어가므로 그대로다. 해는
\[
a^*\approx28{,}993.06,
\qquad b^*\approx86{,}574.07.
\]
두 위험자산 투자액의 합이 100,000달러보다 크므로 채권보유는
\[
100{,}000-a^*-b^*\approx-15{,}567.13.
\]
음의 채권보유는 무위험이자율로 차입하여 위험자산을 추가 매수한다는 뜻이다.

상관계수가 음수가 되면 한 자산이 나쁠 때 다른 자산이 좋을 가능성이 커진다. 두 자산을 함께 더 많이 보유해도 포트폴리오 전체 위험을 상쇄할 수 있으므로 위험자산 수요가 모두 증가하고 레버리지까지 사용한다.
\end{solutionbox}

\begin{answerbox}
\begin{center}
\begin{tabular}{lrrr}
\toprule
경우 & Adidas & Boeing & 무위험채권 \\
\midrule
(a) Adidas+채권 & $22{,}500.00$ & $0$ & $77{,}500.00$ \\
(b) $\rho=0.2$ & $17{,}881.94$ & $61{,}574.07$ & $20{,}543.98$ \\
(c) $\rho=-0.2$ & $28{,}993.06$ & $86{,}574.07$ & $-15{,}567.13$ \\
\bottomrule
\end{tabular}
\end{center}
\end{answerbox}

\begin{warningbox}
분산을 $a^2\sigma_A^2+b^2\sigma_B^2$에서 끝내면 큰 감점이다. 반드시 $2ab\sigma_{AB}$를 먼저 쓰고, 상관계수에서 공분산을 계산해야 한다. 투자액과 포트폴리오 비중을 혼용하지도 말 것.
\end{warningbox}

\exercise{PS3 Exercise 2}{공식 문제}
\begin{officialbox}
$t=0,1$의 두 시점, 무위험수익률 $r_f$, 현금흐름 $\widetilde\delta\sim N(\bar\delta,\sigma_\delta^2)$를 지급하는 위험자산이 있다. 현재가격은 $P$, 총공급은 $S$, 동일한 투자자는 $N$명이다. 각 투자자는 초기부 $Y_0$와 CARA 효용
\[
U(Y)=-e^{-\nu Y},\qquad \nu>0
\]
을 갖는다.
\begin{enumerate}[label=(\alph*)]
\item 투자자가 위험자산 $X$주를 매수할 때 최적화문제를 세워라.
\item 수요함수 $X(P)$를 구하라.
\item 시장청산으로 균형가격 $P$를 구하라.
\item $\sigma_\delta^2$가 증가할 때 가격방향과 직관을 설명하라.
\end{enumerate}
\end{officialbox}

\begin{strategybox}
CARA--Normal에서는 정규분포의 적률생성함수
\[
Z\sim N(\mu,s^2)\quad\Longrightarrow\quad \E[e^Z]=e^{\mu+s^2/2}
\]
를 이용하면 기대효용 최대화가 확실성등가 최대화로 바뀐다.
\end{strategybox}

\begin{solutionbox}
\step{(a) 최종부}
$X$주를 사는 데 $PX$가 들고 나머지를 무위험자산에 둔다.
\[
\widetilde Y_1=\widetilde\delta X+(1+r_f)(Y_0-PX).
\]
따라서 문제는
\[
\max_X\;\E\left[-e^{-\nu\widetilde Y_1}\right].
\]

\step{(b) 확실성등가와 수요}
$-\nu\widetilde Y_1$은 정규분포이며
\[
\E[-\nu\widetilde Y_1]
=-\nu\left[\bar\delta X+(1+r_f)(Y_0-PX)\right],
\]
\[
\Var(-\nu\widetilde Y_1)=\nu^2X^2\sigma_\delta^2.
\]
적률생성함수를 적용하면 기대효용 최대화는 다음 확실성등가 최대화와 동치다.
\[
\max_X\left\{\bar\delta X+(1+r_f)(Y_0-PX)-\frac\nu2X^2\sigma_\delta^2\right\}.
\]
FOC는
\[
\bar\delta-(1+r_f)P-\nu\sigma_\delta^2X=0.
\]
따라서
\[
\boxed{X(P)=\frac{\bar\delta-(1+r_f)P}{\nu\sigma_\delta^2}}.
\]

\step{(c) 시장청산}
총수요 $NX(P)$가 공급 $S$와 같아야 한다.
\[
N\frac{\bar\delta-(1+r_f)P}{\nu\sigma_\delta^2}=S.
\]
따라서
\[
\boxed{P=\frac{\bar\delta-\nu\sigma_\delta^2(S/N)}{1+r_f}}.
\]
이는
\[
P=\frac{\text{기대현금흐름}-\text{위험프리미엄}}{1+r_f}
\]
의 형태다.

\step{(d) 위험 증가}
\[
\frac{\partial P}{\partial\sigma_\delta^2}
=-\frac{\nu S/N}{1+r_f}<0.
\]
현금흐름 변동성이 커지면 위험회피 투자자가 같은 공급량을 보유하기 위해 더 큰 위험프리미엄을 요구한다. 요구수익률이 올라가려면 현재가격은 내려가야 한다.
\end{solutionbox}

\begin{answerbox}
\[
X(P)=\frac{\bar\delta-(1+r_f)P}{\nu\sigma_\delta^2},
\qquad
P=\frac{\bar\delta-\nu\sigma_\delta^2S/N}{1+r_f},
\qquad
\sigma_\delta^2\uparrow\Rightarrow P\downarrow.
\]
\end{answerbox}

\begin{warningbox}
현금흐름 $\widetilde\delta$는 \textbf{가격을 포함한 수익률}이 아니라 $t=1$의 지급액이다. 따라서 최종부에는 $\widetilde\delta X$와 $-PX$가 서로 다른 시점으로 들어간다.
\end{warningbox}

\exercise{PS3 Exercise 3}{공식 문제 · 공식 exercise-session 해설}
\begin{officialbox}
Exercise 2를 수정한다. 투자자의 비율 $\alpha$는 낙관형 H로 $\widetilde\delta\sim N(\bar\delta_H,\sigma^2)$를 믿고, 나머지 $1-\alpha$는 비관형 L로 $\widetilde\delta\sim N(\bar\delta_L,\sigma^2)$를 믿는다. $\bar\delta_L<\bar\delta_H$이고 각 투자자는 시장이 열리기 전 $S/N$주씩 보유한다.
\begin{enumerate}[label=(\alph*)]
\item 경쟁균형 $X_H,X_L,P$를 구하라.
\item 거래량을 구하고, 누가 얼마를 거래하는지와 거래량이 언제 큰지 설명하라.
\end{enumerate}
\end{officialbox}

\begin{strategybox}
유형별 수요함수는 Exercise 2의 평균만 바꾸면 된다. 시장청산에는 유형별 인원수 $\alpha N$과 $(1-\alpha)N$를 곱한다.
\end{strategybox}

\begin{solutionbox}
\step{(a) 유형별 수요와 가격}
\[
X_H(P)=\frac{\bar\delta_H-(1+r_f)P}{\nu\sigma^2},
\qquad
X_L(P)=\frac{\bar\delta_L-(1+r_f)P}{\nu\sigma^2}.
\]
시장청산은
\[
\alpha NX_H(P)+(1-\alpha)NX_L(P)=S.
\]
대입하여
\[
\boxed{P=
\frac{\alpha\bar\delta_H+(1-\alpha)\bar\delta_L-\nu\sigma^2S/N}{1+r_f}}.
\]
가격은 시장 전체의 평균적 기대현금흐름에서 위험프리미엄을 뺀 뒤 할인한 값이다.

가격을 각 수요함수에 다시 넣으면
\[
\boxed{X_H=\frac SN+\frac{(1-\alpha)(\bar\delta_H-\bar\delta_L)}{\nu\sigma^2}},
\]
\[
\boxed{X_L=\frac SN-\frac{\alpha(\bar\delta_H-\bar\delta_L)}{\nu\sigma^2}}.
\]
낙관형은 초기보유량보다 더 사고 비관형은 판다.

\step{(b) 거래량}
낙관형 한 명의 순매수는
\[
X_H-\frac SN=\frac{(1-\alpha)(\bar\delta_H-\bar\delta_L)}{\nu\sigma^2}.
\]
낙관형 전체의 매수량은
\[
\alpha N\left(X_H-\frac SN\right)
=\frac{N\alpha(1-\alpha)(\bar\delta_H-\bar\delta_L)}{\nu\sigma^2}.
\]
비관형 전체의 순매도 절대값도 정확히 같다. 따라서 시장 거래량은
\[
\boxed{V=\frac{N\alpha(1-\alpha)(\bar\delta_H-\bar\delta_L)}{\nu\sigma^2}}.
\]
$\alpha(1-\alpha)$는 역 U자이며 $\alpha=1/2$에서 최대다. 또한 의견차 $\bar\delta_H-\bar\delta_L$가 클수록, 위험회피 $\nu$가 작을수록, 믿음의 분산 $\sigma^2$가 작아 각자 더 확신할수록 거래량이 커진다.
\end{solutionbox}

\begin{answerbox}
\[
P=\frac{\alpha\bar\delta_H+(1-\alpha)\bar\delta_L-\nu\sigma^2S/N}{1+r_f},
\]
\[
X_H=\frac SN+\frac{(1-\alpha)\Delta\bar\delta}{\nu\sigma^2},
\qquad
X_L=\frac SN-\frac{\alpha\Delta\bar\delta}{\nu\sigma^2},
\]
\[
V=\frac{N\alpha(1-\alpha)\Delta\bar\delta}{\nu\sigma^2},\qquad
\Delta\bar\delta=\bar\delta_H-\bar\delta_L>0.
\]
\end{answerbox}

\begin{warningbox}
거래량은 최종보유량 $X_i$ 자체가 아니라 \textbf{초기보유량 $S/N$에서 얼마나 변했는지}로 계산한다. 시장 전체에서 매수량과 매도량을 다시 합쳐 두 배로 세지 않는다.
\end{warningbox}

\exercise{PS3 Exercise 4}{공식 문제 · 공식 exercise-session 해설}
\begin{officialbox}
Exercise 2의 $N$명 가격수용자 외에 가격영향을 이해하는 대형투자자가 한 명 있다. 대형투자자는 위험중립이며 현금흐름 기대값을 $\bar\delta_B\ne\bar\delta$로 믿는다. 대형투자자의 매수량을 $Z$라 한다.
\begin{enumerate}[label=(\alph*)]
\item 경쟁균형 $X,Z,P$를 구하라.
\item 공식 문항이 요구하는 $\partial Z/\partial\sigma<0$의 이유를 설명하라. 위험중립자가 왜 위험이 커질 때 매수를 줄이는가?
\end{enumerate}
\end{officialbox}

\begin{strategybox}
대형투자자가 $Z$주를 사면 가격수용자에게 남는 공급은 $S-Z$다. 먼저 가격수용자의 시장청산에서 $P(Z)$를 구한 뒤, 대형투자자가 그 가격함수를 내생적으로 고려해 $Z$를 선택한다.
\end{strategybox}

\begin{solutionbox}
\step{가격수용자 측 시장청산}
각 소형투자자의 수요는
\[
X(P)=\frac{\bar\delta-(1+r)P}{\nu\sigma^2}.
\]
대형투자자가 $Z$를 보유하면
\[
NX(P)+Z=S.
\]
따라서
\[
\boxed{P(Z)=\frac{\bar\delta-(\nu\sigma^2/N)(S-Z)}{1+r}}.
\]
$Z$가 증가하면 소형투자자에게 남는 공급이 줄고 가격은 상승한다.

\step{대형투자자의 최적화}
대형투자자는 위험중립이므로 기대최종부를 최대화한다.
\[
\max_Z\;\bar\delta_B Z+(1+r)\bigl[Y_0-P(Z)Z\bigr].
\]
$P(Z)$를 대입하고 상수항을 제외하면
\[
\max_Z\left\{
\left(\bar\delta_B-\bar\delta+\frac{\nu\sigma^2S}{N}\right)Z
-\frac{\nu\sigma^2}{N}Z^2
\right\}.
\]
FOC에서
\[
\boxed{Z^*=\frac S2+\frac{N(\bar\delta_B-\bar\delta)}{2\nu\sigma^2}}.
\]
이를 $P(Z)$에 넣으면
\[
\boxed{P^*=\frac1{1+r}\left[
\frac{\bar\delta+\bar\delta_B}{2}-\frac{S\nu\sigma^2}{2N}
\right]}.
\]
소형투자자의 보유량은 시장청산 또는 수요함수로
\[
\boxed{X^*=\frac{S}{2N}-\frac{\bar\delta_B-\bar\delta}{2\nu\sigma^2}}.
\]

\step{(b) 위험과 가격영향}
공식 해설이 설명하는 낙관적 대형투자자 $\bar\delta_B>\bar\delta$의 경우
\[
\frac{\partial Z^*}{\partial\sigma}
=-\frac{N(\bar\delta_B-\bar\delta)}{\nu\sigma^3}<0.
\]
대형투자자가 직접 위험을 싫어해서가 아니다. $\sigma$가 커지면 위험회피 소형투자자의 수요곡선이 더 가팔라져, 대형투자자가 $Z$를 늘려 잔여공급을 줄일 때 가격이 더 크게 상승한다. 자신이 매수하면서 스스로 지불가격을 끌어올리는 \textbf{price impact}가 커지기 때문에 최적매수량을 줄인다.

\smallnote{원문은 $\bar\delta_B\ne\bar\delta$라고 쓰면서 $\partial Z/\partial\sigma<0$을 요구한다. 위 부호는 공식 해설이 의도한 $\bar\delta_B>\bar\delta$에서 성립하며, 반대의 믿음이면 비교정태의 부호도 반대가 된다.}
\end{solutionbox}

\begin{answerbox}
\[
Z^*=\frac S2+\frac{N(\bar\delta_B-\bar\delta)}{2\nu\sigma^2},
\quad
X^*=\frac{S}{2N}-\frac{\bar\delta_B-\bar\delta}{2\nu\sigma^2},
\]
\[
P^*=\frac1{1+r}\left[rac{\bar\delta+\bar\delta_B}{2}-\frac{S\nu\sigma^2}{2N}\right].
\]
낙관적 대형투자자의 경우 $\sigma\uparrow$는 price impact를 키워 $Z^*\downarrow$를 만든다.
\end{answerbox}

\begin{warningbox}
대형투자자를 가격수용자로 처리해 $P$를 상수로 미분하면 이 문제의 핵심이 사라진다. 반드시 $P=P(Z)$를 목적함수 안에 넣어야 한다.
\end{warningbox}

\exercise{PS3 Exercise 5}{공식 문제}
\begin{officialbox}
위험자산은 $t=1$에 $A>0$을 확률 $\pi$로, 0을 확률 $1-\pi$로 지급한다. 현재가격 $P$, 공급 $S$, 투자자 수 $N$, 무위험 순수익률은 $r$다. 각 투자자의 기대효용은
\[
\E[U(\widetilde Y_1)]=\E[\widetilde Y_1]-\frac\theta2\Var(\widetilde Y_1),\qquad\theta>0.
\]
\begin{enumerate}[label=(\alph*)]
\item 수요함수와 균형가격을 구하라.
\item $\theta$가 증가할 때 가격변화를 설명하라.
\item $A$가 증가할 때 가격변화를 설명하라.
\end{enumerate}
\end{officialbox}

\begin{strategybox}
먼저 현금흐름의 평균과 분산을 계산한다. $A$는 평균도 올리지만 분산도 제곱으로 올리므로 가격효과가 단조롭지 않을 수 있다.
\end{strategybox}

\begin{solutionbox}
\step{현금흐름 통계량}
\[
\E[\widetilde\delta]=\pi A,
\qquad
\Var(\widetilde\delta)=\pi(1-\pi)A^2.
\]
최종부는
\[
\widetilde Y_1=\widetilde\delta X+(1+r)(Y_0-PX).
\]
따라서 선택변수와 관련된 목적함수는
\[
\pi AX-(1+r)PX-\frac\theta2\pi(1-\pi)A^2X^2.
\]
FOC에서
\[
\boxed{X(P)=\frac{\pi A-(1+r)P}{\theta\pi(1-\pi)A^2}}.
\]

\step{(a) 균형가격}
$NX(P)=S$를 적용하면
\[
\boxed{P=rac{\pi A}{1+r}
-\frac{S\theta\pi(1-\pi)A^2}{(1+r)N}}.
\]
첫 항은 기대지급액의 할인현재가치, 둘째 항은 위험프리미엄이다.

\step{(b) 위험회피도}
\[
\frac{\partial P}{\partial\theta}
=-\frac{S\pi(1-\pi)A^2}{(1+r)N}<0.
\]
위험회피가 강해질수록 같은 공급을 보유시키려면 더 큰 위험프리미엄과 더 낮은 가격이 필요하다.

\step{(c) 상방지급액 $A$}
\[
\frac{\partial P}{\partial A}
=\frac{\pi}{1+r}\left[1-\frac{2S\theta(1-\pi)A}{N}\right].
\]
따라서
\[
\frac{\partial P}{\partial A}>0
\quad\Longleftrightarrow\quad
A<\frac{N}{2\theta S(1-\pi)}.
\]
작은 $A$에서는 기대현금흐름 증가효과가 우세하여 가격이 상승한다. 큰 $A$에서는 분산과 위험프리미엄이 $A^2$로 커지는 효과가 우세하여 가격이 하락한다. 가격은 $A$에 대해 역 U자다.
\end{solutionbox}

\begin{answerbox}
\[
X(P)=\frac{\pi A-(1+r)P}{\theta\pi(1-\pi)A^2},
\qquad
P=\frac{\pi A}{1+r}-\frac{S\theta\pi(1-\pi)A^2}{(1+r)N}.
\]
$\theta\uparrow\Rightarrow P\downarrow$. $A$에 대해서는 임계값 $N/[2\theta S(1-\pi)]$ 전까지 상승하고 이후 하락한다.
\end{answerbox}

\begin{warningbox}
$A$가 커지면 기대값만 커진다고 결론 내리면 안 된다. 이 분포에서는 분산도 $A^2$로 증가하므로 두 효과를 모두 미분해야 한다.
\end{warningbox}

\exercise{PS3 Exercise 6}{공식 문제}
\begin{officialbox}
위험자산 현금흐름은
\[
\widetilde\delta=
\begin{cases}
\delta_H,&1/2,\\
\delta_L,&1/2,
\end{cases}
\qquad \delta_H>\delta_L.
\]
가격 $P$, 공급 $S$, 투자자 수 $N$, 무위험 순수익률 $r>0$, CARA 계수 $\nu>0$이다.
\begin{enumerate}[label=(\alph*)]
\item 투자자 문제를 세워라.
\item 수요함수 $X(P)$를 구하라.
\item 균형가격을 구하라.
\item $N\to\infty$일 때 가격과 직관을 설명하라.
\end{enumerate}
\end{officialbox}

\begin{strategybox}
이번 현금흐름은 정규분포가 아니므로 CARA--Normal의 확실성등가 공식을 바로 쓰면 안 된다. 두 상태 기대효용을 직접 미분하고 지수항의 비율을 로그로 푼다.
\end{strategybox}

\begin{solutionbox}
\step{(a) 문제}
\[
\widetilde Y_1=\widetilde\delta X+(1+r)(Y_0-PX),
\qquad
\max_X\E[-e^{-\nu\widetilde Y_1}].
\]

\step{(b) FOC와 수요함수}
두 상태를 직접 쓰면 FOC는
\begin{align*}
0={}&\frac12[\delta_H-(1+r)P]
 e^{-\nu[\delta_HX+(1+r)(Y_0-PX)]}\\
&+\frac12[\delta_L-(1+r)P]
 e^{-\nu[\delta_LX+(1+r)(Y_0-PX)]}.
\end{align*}
공통항을 정리하면
\[
\delta_H-(1+r)P
=\bigl[(1+r)P-\delta_L\bigr]e^{\nu(\delta_H-\delta_L)X}.
\]
따라서
\[
\boxed{X(P)=\frac1{\nu(\delta_H-\delta_L)}
\ln\left[\frac{\delta_H-(1+r)P}{(1+r)P-\delta_L}\right]}.
\]
내부해에는 $\delta_L<(1+r)P<\delta_H$가 필요하다.

\step{(c) 시장청산}
$NX(P)=S$에서
\[
\frac{\delta_H-(1+r)P}{(1+r)P-\delta_L}
=\exp\left[\frac{S\nu(\delta_H-\delta_L)}{N}\right].
\]
$K=\exp[S\nu(\delta_H-\delta_L)/N]$라 두면
\[
\delta_H-(1+r)P=K[(1+r)P-\delta_L].
\]
따라서
\[
\boxed{P=\frac1{1+r}\frac{\delta_H+K\delta_L}{1+K}},
\qquad
K=e^{S\nu(\delta_H-\delta_L)/N}.
\]

\step{(d) 투자자 수의 극한}
$N\to\infty$이면 $K\to1$이다.
\[
\lim_{N\to\infty}P
=\frac1{1+r}\frac{\delta_H+\delta_L}{2}
=\frac{\E[\widetilde\delta]}{1+r}.
\]
총공급 $S$가 고정된 상태에서 투자자 수가 늘면 한 사람의 균형보유량 $S/N$이 0으로 간다. 각자가 부담하는 위험이 사라지므로 요구 위험프리미엄도 0으로 수렴한다.
\end{solutionbox}

\begin{answerbox}
\[
X(P)=\frac1{\nu(\delta_H-\delta_L)}
\ln\frac{\delta_H-(1+r)P}{(1+r)P-\delta_L},
\]
\[
P=\frac1{1+r}
\frac{\delta_H+\delta_L e^{S\nu(\delta_H-\delta_L)/N}}
{1+e^{S\nu(\delta_H-\delta_L)/N}},
\]
\[
N\to\infty\quad\Longrightarrow\quad
P\to\frac{\delta_H+\delta_L}{2(1+r)}.
\]
\end{answerbox}

\begin{warningbox}
두 점 분포를 정규분포로 취급해 평균--분산 확실성등가를 쓰면 공식 문항의 해와 달라진다. 여기서는 기대효용을 상태별로 직접 미분한다.
\end{warningbox}

\exercise{PS3 Exercise 7}{공식 문제}
\begin{officialbox}
현금흐름은
\[
\widetilde\delta=
\begin{cases}
\bar\delta+\alpha,&1/2,\\
\bar\delta-\alpha,&1/2,
\end{cases}
\qquad \bar\delta>\alpha>0.
\]
로그효용 $U(Y_1)=\ln Y_1$을 가지며, 단순화를 위해 $r=0$, $S=1$, $N=1$, $Y_0=1$이다.
\begin{enumerate}[label=(\alph*)]
\item 투자자 문제를 세워라.
\item 수요함수 $X(P)$를 구하라.
\item $P\le\bar\delta$인 경우에 초점을 맞춰 균형가격을 구하라.
\item $\alpha$가 증가할 때 가격과 직관을 설명하라.
\end{enumerate}
\end{officialbox}

\begin{strategybox}
두 상태의 최종부를 각각 쓰고 로그 기대효용을 직접 미분한다. 시장청산은 한 명의 투자자가 한 주를 보유하므로 $X(P)=1$이다.
\end{strategybox}

\begin{solutionbox}
\step{(a) 최종부와 목적함수}
$r=0$, $Y_0=1$이므로
\[
\widetilde Y_1=\widetilde\delta X+(1-PX).
\]
따라서
\begin{align*}
\max_X\;&\frac12\ln\bigl[(\bar\delta+\alpha)X+1-PX\bigr]\\
&+\frac12\ln\bigl[(\bar\delta-\alpha)X+1-PX\bigr].
\end{align*}

\step{(b) 수요함수}
FOC는
\[
\frac12\frac{\bar\delta+\alpha-P}{1+(\bar\delta+\alpha-P)X}
+\frac12\frac{\bar\delta-\alpha-P}{1+(\bar\delta-\alpha-P)X}=0.
\]
정리하면
\[
\boxed{X(P)=
\frac{\bar\delta-P}
{[(\bar\delta+\alpha)-P][P-(\bar\delta-\alpha)]}}.
\]

\step{(c) 시장청산}
$X(P)=1$을 적용하면
\[
\bar\delta-P=[(\bar\delta+\alpha)-P][P-(\bar\delta-\alpha)].
\]
전개하면
\[
P^2-(1+2\bar\delta)P+\bar\delta+\bar\delta^2-\alpha^2=0.
\]
근의 공식으로
\[
P=\bar\delta+\frac12\left(1\pm\sqrt{1+4\alpha^2}\right).
\]
힌트대로 $P\le\bar\delta$를 만족하는 작은 근을 선택한다.
\[
\boxed{P=\bar\delta+rac12\left(1-\sqrt{1+4\alpha^2}\right)}.
\]
이를
\[
P=\bar\delta-rac12\left(\sqrt{1+4\alpha^2}-1\right)
\]
로 쓰면 기대지급액에서 위험프리미엄을 뺀 구조가 분명해진다.

\step{(d) 변동폭}
\[
\frac{\partial P}{\partial\alpha}
=-\frac{2\alpha}{\sqrt{1+4\alpha^2}}<0.
\]
$\alpha$는 현금흐름의 표준편차다. 변동폭이 커질수록 로그효용 위험회피자가 요구하는 위험프리미엄이 커져 현재가격이 하락한다.
\end{solutionbox}

\begin{answerbox}
\[
X(P)=
\frac{\bar\delta-P}
{[(\bar\delta+\alpha)-P][P-(\bar\delta-\alpha)]},
\]
\[
P=\bar\delta+\frac12\left(1-\sqrt{1+4\alpha^2}\right),
\qquad
\alpha\uparrow\Rightarrow P\downarrow.
\]
\end{answerbox}

\begin{warningbox}
2차방정식의 두 근을 모두 균형이라고 쓰면 안 된다. 공식 힌트의 $P\le\bar\delta$와 최종부의 양의 조건을 이용해 작은 근을 선택해야 한다.
\end{warningbox}

\chaptertitle{PS4 · 평균--분산 분석과 효율적 프런티어}{공식 Exercise 1--7 · 교수님 5단계 풀이, $\rho=1,0,-1$, MVP, 무위험자산과 공매도}

\begin{sourcebox}
\textbf{이 장의 공식 출처}: MBF\_PS4.pdf, MBF\_PS4\_Solution.pdf 및 공식 exercise session 녹취·수업 필기. 공식 해설 PDF에서 Exercise 1은 ``Discussed in the exercise session''이라고 되어 있으므로 해당 문항만 수업 설명을 근거로 재구성했다.
\end{sourcebox}

\section*{교수님이 강조한 PS4 정석 풀이 5단계}
\addcontentsline{toc}{subsection}{PS4 정석 풀이 5단계}
\begin{strategybox}[title={\bfseries 시험장에서 그대로 쓰는 5단계}]
\begin{enumerate}
\item 각 자산의 $\bar r_A,\bar r_B,\sigma_A,\sigma_B,\sigma_{AB}$를 계산한다.
\item $\rho_{AB}=\sigma_{AB}/(\sigma_A\sigma_B)$를 계산한다. 수업에서 교수님은 시험문제라면 계산이 쉬운 $1,0,-1$이 특히 유력하다고 설명했다.
\item 가중치 $w$를 정의하고
\[
\bar r_p=w\bar r_A+(1-w)\bar r_B,
\]
\[
\sigma_p^2=w^2\sigma_A^2+(1-w)^2\sigma_B^2+2w(1-w)\sigma_{AB}
\]
를 쓴다.
\item 평균식에서 $w$를 풀어 분산식에 대입해 $w$를 제거한다. 단, Exercise 1처럼 $\bar r_A=\bar r_B$이면 이 단계가 $0/0$이 되므로 분산을 직접 최소화한다.
\item $(\sigma_p,\bar r_p)$ 평면에 A와 B를 먼저 찍고, 도출한 식과 공매도 조건을 반영해 frontier를 그린다.
\end{enumerate}
\end{strategybox}

\begin{warningbox}
공분산이 0인 문제도 처음부터 $2w(1-w)\sigma_{AB}$를 빼고 쓰지 않는다. 전체 분산식을 쓴 뒤 ``$\sigma_{AB}=0$이므로 마지막 항은 0''이라고 표시해야 계산을 알고 있다는 점이 드러난다.
\end{warningbox}

\exercise{PS4 Exercise 1}{공식 문제 · 수업 풀이 재구성}
\begin{officialbox}
세 상태가 각각 $1/3$ 확률로 발생한다. 무위험수익률은 $r_f=1$이고 위험자산 수익률은
\[
\begin{array}{c|ccc}
&\text{상태 1}&\text{상태 2}&\text{상태 3}\\\hline
\widetilde r_A&1&2&3\\
\widetilde r_B&3&0&3
\end{array}
\]
이다.
\begin{enumerate}[label=(\alph*)]
\item A와 B만 이용하고 두 자산의 공매도가 금지될 때 효율적 프런티어를 구하고 그려라. 분산을 최소화하는 가중치를 찾으라는 힌트가 주어져 있다.
\item A, B와 무위험자산을 이용하며 $r_f$로 자유롭게 차입·대출할 수 있다. A와 B의 공매도는 금지된다. 효율적 프런티어와 자산배분을 설명하라.
\item ``위험회피도가 높을수록 B의 비중은 낮고 A의 비중은 높다''는 진술은 참, 거짓, 판단불가 중 무엇인가?
\item ``금리인하는 Sharpe ratio를 높이므로 모든 투자자에게 환영받는다''는 진술은 참, 거짓, 판단불가 중 무엇인가?
\end{enumerate}
\end{officialbox}

\begin{strategybox}
이 문항은 교수님이 정석 5단계의 \textbf{예외}라고 강조했다. 두 위험자산의 기대수익률이 모두 2라서 평균식으로 $w$를 풀 수 없다. 따라서 포트폴리오 분산을 $w$의 함수로 직접 최소화한다.
\end{strategybox}

\begin{solutionbox}
\step{1--2단계: 통계량}
\[
\bar r_A=\frac{1+2+3}{3}=2,
\qquad
\bar r_B=\frac{3+0+3}{3}=2.
\]
\[
\sigma_A^2=\frac{(1-2)^2+(2-2)^2+(3-2)^2}{3}=\frac23,
\]
\[
\sigma_B^2=\frac{(3-2)^2+(0-2)^2+(3-2)^2}{3}=2.
\]
\[
\sigma_{AB}=\frac{(1-2)(3-2)+(2-2)(0-2)+(3-2)(3-2)}3=0.
\]
따라서 $\rho_{AB}=0$이다.

\step{(a) 위험자산만 있을 때}
A의 비중을 $w$, B의 비중을 $1-w$라 두고 공매도 금지이므로 $0\le w\le1$이다.
\[
\bar r_p=w(2)+(1-w)(2)=2.
\]
모든 포트폴리오의 기대수익률이 동일하다. 분산은
\[
\sigma_p^2=w^2\frac23+(1-w)^2(2)+2w(1-w)(0).
\]
전개하거나 완전제곱하면
\[
\sigma_p^2=\frac83w^2-4w+2
=\frac83\left(w-\frac34\right)^2+\frac12.
\]
따라서 최소분산 가중치는
\[
w_{MVP}=\frac34,\qquad1-w_{MVP}=\frac14,
\]
\[
\sigma_{MVP}^2=\frac12,\qquad \sigma_{MVP}=\frac1{\sqrt2}.
\]
다른 모든 포트폴리오는 기대수익률은 2로 같으면서 위험만 더 크므로 MVP에 의해 지배된다. 따라서 효율적 프런티어는 곡선 전체가 아니라
\[
\boxed{\left(\sigma_p,\bar r_p\right)=\left(\frac1{\sqrt2},2\right)}
\]
한 점이다.

\step{(b) 무위험자산 추가}
위험자산 기회집합 중 같은 기대수익률 2에서 위험이 가장 작은 MVP가 접점포트폴리오 $T$가 된다. $T$의 위험자산 내부 구성은
\[
A:B=\frac34:\frac14=3:1.
\]
무위험점 $F=(0,1)$과 $T=(1/\sqrt2,2)$를 잇는 자본배분선의 기울기는
\[
\frac{2-1}{1/\sqrt2}=\sqrt2.
\]
차입과 대출이 모두 가능하므로 효율적 프런티어는
\[
\boxed{\bar r_p=1+\sqrt2\,\sigma_p,\qquad\sigma_p\ge0}.
\]
투자자가 총부의 비율 $y$를 $T$에 넣으면
\[
\sigma_p=y\sigma_T,\qquad \bar r_p=r_f+y(\bar r_T-r_f).
\]
$0<y<1$이면 무위험자산을 양으로 보유하는 순대출자, $y=1$이면 $T$만 보유, $y>1$이면 무위험자산을 음으로 보유하는 차입자다. 어느 경우에도 위험자산 내부 A:B 비율은 3:1로 같다.

\step{(c) 2기금 분리}
진술은 거짓이다. 위험회피도는 $T$와 무위험자산 사이의 비율 $y$를 바꾼다. 그러나 모든 투자자가 선택하는 위험자산 바스켓 $T$의 내부비율은 A:B=3:1로 동일하다. 더 위험선호적인 투자자가 B만 더 늘리는 것이 아니다.

\step{(d) 금리인하의 후생효과}
위험자산 $T$의 기대수익률과 위험이 고정되어 있다고 하자. 같은 $y$에서
\[
\bar r_p=r_f+y(\bar r_T-r_f)
=y\bar r_T+(1-y)r_f.
\]
따라서
\[
\frac{\partial\bar r_p}{\partial r_f}=1-y.
\]
$r_f$가 하락하면
\begin{itemize}
\item $y<1$인 순대출자는 안전자산 수익이 낮아져 불리하다.
\item $y=1$인 $T$ 단독보유자는 직접 영향이 없다.
\item $y>1$인 순차입자는 차입비용이 낮아져 유리하다.
\end{itemize}
그러므로 ``모든 투자자''가 좋아한다는 결론은 낼 수 없고 답은 판단불가다.
\end{solutionbox}

\begin{answerbox}
\begin{align*}
\text{(a)}\quad &w_A=\frac34,\;w_B=\frac14,\;\sigma_{MVP}=\frac1{\sqrt2},\;\bar r_{MVP}=2,\\
&\text{효율적 프런티어는 MVP 한 점.}\\
\text{(b)}\quad &\bar r_p=1+\sqrt2\sigma_p,\;\sigma_p\ge0,\quad T:\;A:B=3:1.\\
\text{(c)}\quad &\text{거짓. 위험회피도는 $T$와 무위험자산의 비율만 바꾼다.}\\
\text{(d)}\quad &\text{판단불가. 대출자는 불리하고 차입자는 유리하다.}
\end{align*}
\end{answerbox}

\begin{warningbox}
$\bar r_A=\bar r_B$인데
\[
w=\frac{\bar r_B-\bar r_p}{\bar r_B-\bar r_A}
\]
를 쓰면 분모가 0이다. 이 문항에서는 반드시 분산을 직접 최소화해야 한다. 또한 효율적 프런티어를 A--B 전체 수평선으로 그리면 같은 수익률에서 더 위험한 점까지 포함하므로 오답이다.
\end{warningbox}

\exercise{PS4 Exercise 2}{공식 문제}
\begin{officialbox}
두 상태가 각각 $1/2$ 확률로 발생하고 무위험자산은 없다. 공매도는 허용된다.
\[
\begin{array}{c|cc}
&\text{상태 1}&\text{상태 2}\\\hline
\widetilde r_A&0&0.2\\
\widetilde r_B&0.8&0
\end{array}
\]
\begin{enumerate}[label=(\alph*)]
\item A와 B 포트폴리오의 효율적 프런티어 식과 그림을 구하라.
\item 초기부 $Y_0=1$이고 $\E[U(\widetilde Y_1)]=\E(\widetilde Y_1)-\frac12\Var(\widetilde Y_1)$인 투자자의 최적 가중치를 구하라.
\item 공매도가 금지되면 효율적 프런티어와 투자자 후생은 어떻게 변하는가?
\end{enumerate}
\end{officialbox}

\begin{strategybox}
상관계수가 $-1$이므로 분산은 완전제곱이 된다. 효율적 프런티어는 위험 0 포트폴리오에서 시작하는 직선이며, 공매도가 허용되면 B를 넘어 북동쪽으로 연장된다.
\end{strategybox}

\begin{solutionbox}
\step{통계량}
\[
\bar r_A=0.1,\quad \bar r_B=0.4,
\quad\sigma_A=0.1,\quad\sigma_B=0.4,
\]
\[
\sigma_{AB}=\frac{(0-0.1)(0.8-0.4)+(0.2-0.1)(0-0.4)}2=-0.04,
\]
\[
\rho_{AB}=\frac{-0.04}{0.1\cdot0.4}=-1.
\]

\step{(a) 프런티어}
A 비중을 $w$라 하면
\[
\bar r_p=0.1w+0.4(1-w)=0.4-0.3w,
\]
\[
\sigma_p^2=[0.1w-0.4(1-w)]^2=(0.5w-0.4)^2.
\]
따라서
\[
\pm\sigma_p=0.5w-0.4
\quad\Longrightarrow\quad
w=0.8\pm2\sigma_p.
\]
평균식에 넣으면 가능한 포트폴리오의 두 직선은
\[
\bar r_p=\frac4{25}\pm\frac35\sigma_p.
\]
같은 위험에서 기대수익률이 높은 위쪽만 효율적이므로
\[
\boxed{\bar r_p=\frac4{25}+\frac35\sigma_p,\qquad\sigma_p\ge0}.
\]
위험 0 포트폴리오는 $w=0.8$, 즉 A 80\%, B 20\%이며 기대수익률은 $0.16$이다.

\step{(b) 최적 가중치}
최종부는 $(1+\widetilde r_p)Y_0$이고 $Y_0=1$이므로 목적함수는
\[
1+(0.4-0.3w)-\frac12(0.5w-0.4)^2.
\]
FOC는
\[
-0.3-0.5(0.5w-0.4)=0.
\]
따라서
\[
\boxed{w_A^*=-\frac25,\qquad w_B^*=1-w_A^*=\frac75}.
\]
A를 40\% 공매도하고 그 자금까지 합쳐 B를 140\% 보유한다.

\step{(c) 공매도 금지}
$0\le w\le1$이므로 $w<0$에 해당하는 B 북동쪽 연장선이 제거된다. 무제약 최적점 $w=-2/5$는 불가능하므로 제약하에서 가장 가까운 효율적 경계 $w=0$, 즉 B 100\%가 최적이다. 선택집합이 줄어들어 최대 기대효용은 낮아진다.
\end{solutionbox}

\begin{answerbox}
\[
\bar r_p=\frac4{25}+\frac35\sigma_p,\quad\sigma_p\ge0,
\qquad
(w_A^*,w_B^*)=\left(-\frac25,\frac75\right).
\]
공매도 금지 시 효율적 프런티어는 위험 0 포트폴리오부터 B까지만 남고, 투자자는 B 100\%를 선택하며 더 나빠진다.
\end{answerbox}

\begin{warningbox}
$\rho=-1$이라고 해서 모든 가중치에서 위험이 0인 것은 아니다. $w\sigma_A=(1-w)\sigma_B$를 만족하는 단 하나의 가중치에서만 위험이 완전히 상쇄된다.
\end{warningbox}

\exercise{PS4 Exercise 3}{공식 문제}
\begin{officialbox}
세 상태가 각각 $1/3$ 확률이며 무위험자산은 없다.
\[
\begin{array}{c|ccc}
&1&2&3\\\hline
\widetilde r_A&1&2&0\\
\widetilde r_B&0&3&3
\end{array}
\]
\begin{enumerate}[label=(\alph*)]
\item A와 B 사이에 평균--분산 지배관계가 있는가?
\item 두 자산 모두 공매도가 허용될 때 minimum-variance frontier를 구하고 그려라. MVP의 평균과 표준편차를 구하면 보너스다.
\item 초기부 1이고 $\E[U(\widetilde Y_1)]=\E(\widetilde Y_1)-\frac16\Var(\widetilde Y_1)$인 투자자의 최적 가중치를 구하고 무차별곡선을 설명하라.
\end{enumerate}
\end{officialbox}

\begin{strategybox}
공분산이 0인 전형적 bullet-shaped frontier다. 평균식에서 $w=2-\bar r_p$를 구해 분산식에 넣는다. MVP와 투자자 최적점은 서로 다른 문제다.
\end{strategybox}

\begin{solutionbox}
\step{통계량과 (a)}
\[
\bar r_A=1,\quad\bar r_B=2,
\quad\sigma_A^2=\frac23,\quad\sigma_B^2=2,
\quad\sigma_{AB}=0.
\]
B는 기대수익률이 높지만 위험도 높다. 따라서 어느 자산도 다른 자산을 평균--분산 지배하지 않는다.

\step{(b) minimum-variance frontier}
A 비중을 $w$라 하면
\[
\bar r_p=w(1)+(1-w)(2)=2-w
\quad\Longrightarrow\quad w=2-\bar r_p.
\]
분산은
\[
\sigma_p^2=w^2\frac23+(1-w)^2(2)+2w(1-w)(0).
\]
$w=2-\bar r_p$, $1-w=\bar r_p-1$을 넣으면
\[
\boxed{\sigma_p=\sqrt{\frac23(2-\bar r_p)^2+2(\bar r_p-1)^2}}.
\]
MVP는 분산을 $w$로 미분하여
\[
\frac{\dd\sigma_p^2}{\dd w}=\frac43w-4(1-w)=0
\quad\Longrightarrow\quad w_{MVP}=\frac34.
\]
따라서
\[
\bar r_{MVP}=2-\frac34=\frac54,
\qquad
\sigma_{MVP}^2=\frac12,
\qquad
\sigma_{MVP}=\frac1{\sqrt2}.
\]

\step{(c) 투자자 최적점}
$Y_0=1$이므로
\[
\E[U]=1+\bar r_p-\frac16\sigma_p^2.
\]
$w$로 쓰면
\[
\E[U]=1+[2-w]-\frac16\left[\frac23w^2+2(1-w)^2\right].
\]
FOC는
\[
-1-\frac16\left(\frac{16}{3}w-4\right)=0
\quad\Longrightarrow\quad
\boxed{w_A^*=-\frac38,\qquad w_B^*=\frac{11}{8}}.
\]
A를 공매도하고 B를 레버리지 보유한다. 참고로
\[
\bar r_p^*=2-\left(-\frac38\right)=\frac{19}{8},
\]
\[
(\sigma_p^*)^2=\frac23\left(\frac{9}{64}\right)+2\left(\frac{11}{8}\right)^2
=\frac{31}{8}.
\]
무차별곡선은 일정 효용 $\overline U$에 대해
\[
\bar r_p=\overline U-1+\frac16\sigma_p^2
\]
이므로 위로 볼록하며 효율적 상단 branch에 접한다.
\end{solutionbox}

\begin{answerbox}
\begin{align*}
\text{(a)}\quad &\text{지배관계 없음.}\\
\text{(b)}\quad &\sigma_p=\sqrt{\frac23(2-\bar r_p)^2+2(\bar r_p-1)^2},\\
&w_{MVP}=\frac34,\quad \bar r_{MVP}=\frac54,\quad\sigma_{MVP}=\frac1{\sqrt2}.\\
\text{(c)}\quad &(w_A^*,w_B^*)=\left(-\frac38,\frac{11}{8}\right).
\end{align*}
\end{answerbox}

\begin{warningbox}
minimum-variance frontier 전체와 efficient frontier 상단을 구분해야 한다. MVP의 왼쪽·아래쪽 branch는 같은 위험에서 더 높은 수익률을 주는 상단 branch에 지배될 수 있다.
\end{warningbox}

\exercise{PS4 Exercise 4}{공식 문제}
\begin{officialbox}
세 상태가 각각 $1/3$ 확률이고 무위험수익률은 $r_f=5.5$다.
\[
\begin{array}{c|ccc}
&1&2&3\\\hline
\widetilde r_A&8&6&4\\
\widetilde r_B&11&7&3
\end{array}
\]
\begin{enumerate}[label=(\alph*)]
\item 평균--분산 지배관계가 있는가?
\item Sharpe ratio로 A와 B를 비교하라.
\item A와 B만 이용하고 공매도가 금지될 때 효율적 프런티어를 구하라.
\item 무위험자산까지 이용하고 자유롭게 차입·대출할 때 자산배분을 설명하라.
\item $r_f$가 4.5로 하락하면 (d)가 어떻게 바뀌며 투자자 후생은 어떻게 되는가?
\end{enumerate}
\end{officialbox}

\begin{strategybox}
두 위험자산은 $\rho=1$이다. 따라서 A--B 기회집합은 직선이고 분산효과가 없다. 무위험자산이 추가되면 두 자산 중 Sharpe ratio가 더 큰 자산과 무위험자산의 조합이 선택된다.
\end{strategybox}

\begin{solutionbox}
\step{통계량}
\[
\bar r_A=6,\quad\bar r_B=7,
\quad\sigma_A^2=\frac83,\quad\sigma_B^2=\frac{32}{3},
\quad\sigma_{AB}=\frac{16}{3}.
\]
$\sigma_B=2\sigma_A$이고
\[
\rho_{AB}=\frac{16/3}{\sqrt{8/3}\sqrt{32/3}}=1.
\]

\step{(a) 지배관계}
B는 평균이 높지만 분산도 높다. 따라서 지배관계는 없다.

\step{(b) Sharpe ratio, $r_f=5.5$}
\[
SR_A=\frac{6-5.5}{\sigma_A}=\frac{0.5}{\sigma_A},
\]
\[
SR_B=\frac{7-5.5}{\sigma_B}=\frac{1.5}{2\sigma_A}=\frac{0.75}{\sigma_A}.
\]
따라서 $SR_B>SR_A$다.

\step{(c) 위험자산만}
A 비중을 $w\in[0,1]$라 하면 $\rho=1$이므로
\[
\sigma_p=w\sigma_A+(1-w)\sigma_B,
\qquad
\bar r_p=6w+7(1-w).
\]
$\sigma_B=2\sigma_A$를 이용해 $w=2-\sigma_p/\sigma_A$를 얻고
\[
\bar r_p=5+\frac{\sigma_p}{\sigma_A}
=5+\sqrt{\frac38}\sigma_p.
\]
따라서
\[
\boxed{\bar r_p=5+\sqrt{\frac38}\sigma_p,\qquad
\sigma_p\in[\sigma_A,\sigma_B]}.
\]

\step{(d) 무위험자산 추가}
$r_f=5.5$에서는 B의 Sharpe ratio가 더 높다. 따라서 투자자는 A를 보유하지 않고 무위험자산과 B를 조합한다. 위험회피도가 높은 사람은 무위험자산 비중이 높고, 위험선호가 강한 사람은 무위험자산을 음으로 보유해 차입한 뒤 B를 100\% 이상 보유할 수 있다.

\step{(e) $r_f=4.5$}
새 Sharpe ratio는
\[
SR_A'=\frac{6-4.5}{\sigma_A}=\frac{1.5}{\sigma_A},
\]
\[
SR_B'=\frac{7-4.5}{2\sigma_A}=\frac{1.25}{\sigma_A}.
\]
이제 $SR_A'>SR_B'$이므로 접점자산이 B에서 A로 바뀐다. 투자자들은 무위험자산과 A를 조합한다.

금리인하의 후생효과는 일률적이지 않다. 무위험자산을 양으로 보유하는 순대출자는 이자수입 하락으로 불리할 수 있고, 차입자는 자금조달비용 하락으로 유리할 수 있다. 선호와 기존 포지션에 따라 더 좋아질 수도, 나빠질 수도 있다.
\end{solutionbox}

\begin{answerbox}
\begin{enumerate}[label=(\alph*)]
\item 평균--분산 지배 없음.
\item $r_f=5.5$에서 $SR_B>SR_A$.
\item $\bar r_p=5+\sqrt{3/8}\,\sigma_p$, $\sigma_p\in[\sigma_A,\sigma_B]$.
\item 무위험자산과 B만 조합한다.
\item $r_f=4.5$에서는 무위험자산과 A만 조합한다. 후생효과는 투자자에 따라 다르다.
\end{enumerate}
\end{answerbox}

\begin{warningbox}
Sharpe ratio는 무조건 $\bar r/\sigma$가 아니라 \textbf{$(\bar r-r_f)/\sigma$}다. 금리가 바뀌면 순위도 바뀔 수 있다.
\end{warningbox}

\exercise{PS4 Exercise 5}{공식 문제}
\begin{officialbox}
세 상태가 각각 $1/3$ 확률이고 무위험자산은 없다.
\[
\begin{array}{c|ccc}
&1&2&3\\\hline
\widetilde r_A&6&8&4\\
\widetilde r_B&7&3&11
\end{array}
\]
\begin{enumerate}[label=(\alph*)]
\item 평균--분산 지배관계가 있는가?
\item 두 자산 모두 공매도 금지일 때 효율적 프런티어를 구하라.
\item B의 공매도만 허용하면 프런티어와 후생은 어떻게 변하는가?
\item A와 B 모두 공매도를 허용하면 어떻게 변하는가?
\end{enumerate}
\end{officialbox}

\begin{strategybox}
$\rho=-1$이고 B가 고수익·고위험자산이다. 위험 0 포트폴리오에서 B 방향이 효율적 상단이고, 반대 방향인 B 공매도는 지배되는 하단을 늘릴 뿐이다.
\end{strategybox}

\begin{solutionbox}
\step{통계량과 (a)}
\[
\bar r_A=6,\quad\bar r_B=7,
\quad\sigma_A^2=\frac83,\quad\sigma_B^2=\frac{32}{3},
\quad\sigma_{AB}=-\frac{16}{3},\quad\rho=-1.
\]
B는 기대수익률과 위험이 모두 높으므로 지배관계가 없다.

\step{(b) 공매도 금지}
A 비중을 $w\in[0,1]$라 하면
\[
\sigma_p=\left|w\sigma_A-(1-w)\sigma_B\right|,
\qquad
\bar r_p=6w+7(1-w)=7-w.
\]
위험 0 가중치는
\[
w_F=\frac{\sigma_B}{\sigma_A+\sigma_B}=\frac23,
\]
그때 기대수익률은
\[
\bar r_F=6\cdot\frac23+7\cdot\frac13=\frac{19}{3}.
\]
위험 0점에서 B로 가는 상단의 기울기는
\[
\frac{7-19/3}{\sigma_B-0}=\frac{1/3}{\sigma_B}=\frac1{\sqrt{24}}.
\]
따라서
\[
\boxed{\bar r_p=\frac{19}{3}+\sqrt{\frac1{24}}\sigma_p,\qquad
\sigma_p\in[0,\sigma_B]}.
\]

\step{(c) B만 공매도 허용}
B 비중 $1-w<0$, 즉 $w>1$인 영역이 추가된다. 이 방향은 A를 지나 기대수익률이 더 낮아지고 위험은 더 커지는 하단 branch다. 이미 존재하는 효율적 포트폴리오에 지배되므로 \textbf{효율적 프런티어는 바뀌지 않는다}. 투자자의 최적 선택과 기대효용도 변하지 않는다.

\step{(d) 두 자산 모두 공매도 허용}
A 공매도는 $w<0$을 허용하여 B를 넘어 더 높은 기대수익률과 위험의 북동쪽 상단을 추가한다. 따라서
\[
\boxed{\bar r_p=\frac{19}{3}+\sqrt{\frac1{24}}\sigma_p,\qquad
\sigma_p\in[0,\infty)}.
\]
선택집합이 확대되므로 어떤 투자자도 나빠지지 않는다. 원래 최적점에 머무는 투자자는 무차별이고, A 공매도를 원하는 투자자는 엄격히 좋아진다.
\end{solutionbox}

\begin{answerbox}
지배관계는 없다. 공매도 금지 frontier는
\[
\bar r_p=\frac{19}{3}+\sqrt{\frac1{24}}\sigma_p,\quad0\le\sigma_p\le\sigma_B.
\]
B 공매도만 허용해도 효율적 프런티어와 후생은 변하지 않는다. 두 자산 모두 허용하면 같은 직선이 $\sigma_p\to\infty$까지 연장되고 투자자는 약하게 더 좋아진다.
\end{answerbox}

\begin{warningbox}
``공매도 허용 = 효율적 프런티어가 항상 확대''는 아니다. 어느 자산을 공매도하는 방향이 상단인지 하단인지, 즉 기대수익률과 위험이 함께 어떻게 변하는지를 확인해야 한다.
\end{warningbox}

\exercise{PS4 Exercise 6}{공식 문제}
\begin{officialbox}
세 상태가 각각 $1/3$ 확률이고 무위험자산은 없으며 공매도가 금지된다.
\[
\begin{array}{c|ccc}
&1&2&3\\\hline
\widetilde r_A&2&4&0\\
\widetilde r_B&4&0&8
\end{array}
\]
\begin{enumerate}[label=(\alph*)]
\item 효율적 프런티어를 구하라.
\item ``극도로 위험회피적인 투자자는 B보다 변동성이 작은 A에 70\% 넘게 투자할 수 있다''는 진술을 판정하라.
\item B를 상태별로 지배하는 새 자산 $C=(4,6,8)$로 교체한다. 평균--분산 투자자가 더 좋아질 수도 나빠질 수도 있음을 그림으로 설명하고, 상태별 지배에도 이러한 역설이 생기는 이유를 설명하라.
\end{enumerate}
\end{officialbox}

\begin{strategybox}
(a)는 Exercise 5와 같은 $\rho=-1$ 구조다. (b)는 낮은 개별 변동성이 아니라 \textbf{포트폴리오 지배}를 봐야 한다. (c)는 평균--분산 기준의 한계를 묻는다.
\end{strategybox}

\begin{solutionbox}
\step{(a) 통계량과 프런티어}
\[
\bar r_A=2,\quad\bar r_B=4,
\quad\sigma_A^2=\frac83,\quad\sigma_B^2=\frac{32}{3},
\quad\sigma_{AB}=-\frac{16}{3},\quad\rho=-1.
\]
A 비중을 $w\in[0,1]$라 하면 위험 0 가중치는
\[
w_F=\frac{\sigma_B}{\sigma_A+\sigma_B}=\frac23,
\]
그때 기대수익률은
\[
\bar r_F=2\cdot\frac23+4\cdot\frac13=\frac83.
\]
상단 기울기는
\[
\frac{4-8/3}{\sigma_B}=\frac{4/3}{\sqrt{32/3}}=\frac1{\sqrt6}.
\]
따라서
\[
\boxed{\bar r_p=\frac83+\sqrt{\frac16}\sigma_p,\qquad
0\le\sigma_p\le\sigma_B}.
\]

\step{(b) 70\% A 진술}
위험 0 포트폴리오의 A 비중이 정확히 $2/3\approx66.7\%$다. $w>2/3$으로 A를 더 늘리면 위험은 다시 커지지만 기대수익률은 낮아진다. 이는 효율적 상단이 아니라 지배되는 하단이다. 따라서 어떤 평균--분산 투자자도 $w_A\ge0.7>2/3$을 선택하지 않는다. 진술은 거짓이다.

\step{(c) B를 C로 교체}
$C=(4,6,8)$은 모든 상태에서 B보다 같거나 높은 수익을 준다. 그러나 A와 C의 통계량은
\[
\bar r_C=6,\quad\sigma_C^2=\frac83,
\quad\sigma_{AC}=-\frac43,\quad\rho_{AC}=-\frac12.
\]
A--C 조합은 평균수익률을 크게 올리지만 $\rho=-1$이 아니므로 위험을 완전히 0으로 만들 수 없다. 높은 평균을 중시하는 투자자는 새 기회집합에서 더 좋아질 수 있고, 위험 0 근처를 강하게 선호하는 투자자는 평균--분산 그림상 더 나빠질 수 있다.

이 결론은 상태별 지배와 충돌하는 것처럼 보이지만, 원인은 Modern Portfolio Theory가 분포 전체가 아니라 평균과 분산만으로 선호를 표현한다는 강한 가정이다. 실제로 C의 상태 2 초과수익 6을 그냥 버리면 B를 복제할 수 있으므로, 완전한 상태별 선택이 허용되는 합리적 투자자는 C 도입으로 나빠질 이유가 없다. 평균--분산 기준이 그 우월성을 완전히 포착하지 못한 것이다.
\end{solutionbox}

\begin{answerbox}
\[
\bar r_p=\frac83+\sqrt{\frac16}\sigma_p,\quad0\le\sigma_p\le\sigma_B.
\]
(b)는 거짓이며 효율적인 A 비중의 최대는 위험 0점의 $2/3$이다. (c)의 역설은 상태별 지배를 평균과 분산 두 숫자만으로 완전히 표현하지 못하는 MPT의 한계에서 나온다.
\end{answerbox}

\begin{warningbox}
개별자산 A가 덜 위험하다는 사실만 보고 A 비중을 무한히 늘리면 안 된다. 포트폴리오 위험은 공분산과 조합비율에 의해 결정되며, 위험 0점을 지나면 A를 더 늘리는 것이 오히려 열등해진다.
\end{warningbox}

\exercise{PS4 Exercise 7}{공식 문제}
\begin{officialbox}
세 상태가 각각 $1/3$ 확률이고 무위험자산은 없으며 공매도가 허용된다.
\[
\begin{array}{c|ccc}
&1&2&3\\\hline
\widetilde r_A&2&5&2\\
\widetilde r_B&0&6&0
\end{array}
\]
투자자의 초기부는 $Y_0$이고 $\widetilde Y_1=(1+\widetilde r_p)Y_0$, $\widetilde r_p=w\widetilde r_A+(1-w)\widetilde r_B$, 효용은 $\ln\widetilde Y_1$이다.
\begin{enumerate}[label=(\alph*)]
\item 효율적 프런티어를 구하라.
\item MVP의 A와 B 가중치를 구하라.
\item 투자자의 최적 가중치, 최적 기대수익률과 표준편차를 구하라.
\end{enumerate}
\end{officialbox}

\begin{strategybox}
$\rho=1$이지만 공매도를 허용하므로 두 자산의 변동을 반대로 배치해 위험 0 포트폴리오를 만들 수 있다. 로그효용은 세 상태 최종부를 직접 넣어 $w$로 미분한다.
\end{strategybox}

\begin{solutionbox}
\step{통계량}
\[
\bar r_A=3,\quad\bar r_B=2,
\quad\sigma_A^2=2,\quad\sigma_B^2=8,
\quad\sigma_{AB}=4,
\]
\[
\rho_{AB}=\frac4{\sqrt2\sqrt8}=1.
\]
A는 B보다 평균이 높고 분산이 낮으므로 A가 B를 평균--분산 지배한다.

\step{(a) 효율적 프런티어}
A 비중을 $w$라 하면
\[
\bar r_p=3w+2(1-w)=2+w.
\]
$\rho=1$이므로
\[
\sigma_p^2=[w\sigma_A+(1-w)\sigma_B]^2
=[w\sqrt2+(1-w)2\sqrt2]^2
=2(2-w)^2.
\]
따라서
\[
\sigma_p=\sqrt2|2-w|.
\]
$w=2$에서 위험이 0이고 평균은 4다. $w\ge2$인 상단 branch에서는
\[
w=2+\frac{\sigma_p}{\sqrt2}.
\]
평균식에 넣으면 효율적 프런티어는
\[
\boxed{\bar r_p=4+\frac1{\sqrt2}\sigma_p,\qquad\sigma_p\ge0}.
\]
다른 branch $\bar r_p=4-\sigma_p/\sqrt2$는 같은 위험에서 더 낮은 평균이므로 지배된다.

\step{(b) MVP}
\[
\sigma_p=0\quad\Longleftrightarrow\quad w=2.
\]
따라서
\[
\boxed{w_A^{MVP}=2,\qquad w_B^{MVP}=1-w=-1}.
\]
B를 100\% 공매도하고 A를 200\% 보유하면 두 자산의 완전한 양의 상관 변동이 상쇄된다.

\step{(c) 로그효용 최적화}
상태별 포트폴리오 수익률은
\[
\widetilde r_{p,1}=2w,\qquad
\widetilde r_{p,2}=5w+6(1-w)=6-w,\qquad
\widetilde r_{p,3}=2w.
\]
따라서 $Y_0$의 로그는 상수이므로
\[
\max_w\left\{\frac23\ln(1+2w)+\frac13\ln(7-w)\right\},
\]
정의역은 $-1/2<w<7$이다. FOC는
\[
\frac{4}{3(1+2w)}-\frac1{3(7-w)}=0.
\]
따라서
\[
4(7-w)=1+2w
\quad\Longrightarrow\quad
\boxed{w_A^*=\frac92,\qquad w_B^*=-\frac72}.
\]
최적 기대수익률은
\[
\bar r_p^*=2+\frac92=\frac{13}{2}.
\]
표준편차는
\[
\sigma_p^*=\sqrt2\left|2-\frac92\right|
=\boxed{\frac{5\sqrt2}{2}}.
\]
이는 효율적 프런티어와 로그효용 투자자의 무차별곡선이 접하는 점이다.
\end{solutionbox}

\begin{answerbox}
\begin{align*}
\text{(a)}\quad &\bar r_p=4+\frac{\sigma_p}{\sqrt2},\quad\sigma_p\ge0.\\
\text{(b)}\quad &(w_A^{MVP},w_B^{MVP})=(2,-1),\quad(\sigma,\bar r)=(0,4).\\
\text{(c)}\quad &(w_A^*,w_B^*)=\left(\frac92,-\frac72\right),\\
&\bar r_p^*=\frac{13}{2},\qquad\sigma_p^*=\frac{5\sqrt2}{2}.
\end{align*}
\end{answerbox}

\begin{warningbox}
$\rho=1$이면 ``분산효과가 전혀 없다''는 설명은 가중치가 $0$과 1 사이일 때의 통상적 long-only 해석이다. 공매도가 허용되면 한 자산을 음으로 보유하여 완전 양의 상관 변동도 상쇄할 수 있다.
\end{warningbox}


\clearpage
\section*{최종 체크리스트}
\addcontentsline{toc}{section}{최종 체크리스트}
\begin{tcolorbox}[colback=mint,colframe=green,boxrule=0.8pt,arc=2mm]
\begin{itemize}
\item PS1: 개인별 수요함수와 시장청산으로 가격비율을 구할 수 있는가?
\item PS1: 계약곡선과 경쟁균형을 ``효율성''과 ``공정성''으로 구분해 설명할 수 있는가?
\item PS2: ARA/RRA 조건만 보고 CARA/CRRA 효용함수를 복원할 수 있는가?
\item PS2: 확실성등가와 위험프리미엄을 기대효용 등식으로 계산할 수 있는가?
\item PS3: 최종부를 쓴 뒤 기대값·분산을 구하고 FOC를 세울 수 있는가?
\item PS3: 개인 수요를 합산해 시장청산으로 균형가격을 구할 수 있는가?
\item PS4: 평균·표준편차·공분산·상관계수를 계산한 뒤 가중치를 제거할 수 있는가?
\item PS4: $\rho=1,0,-1$ 각각의 frontier 모양과 공매도 제약을 설명할 수 있는가?
\item PS4: 무위험자산이 있을 때 접점포트폴리오, 대출자·차입자, 2기금 분리를 설명할 수 있는가?
\end{itemize}
\end{tcolorbox}

\end{document}
